\documentclass[reprint,twocolumn,superscriptaddress,nofootinbib]{revtex4-2}

\usepackage{graphicx}
\usepackage{amsmath,amssymb,amsthm,mathrsfs,amsfonts}
\usepackage{mathtools}
\usepackage{bm}
\usepackage{xcolor}
\usepackage{hyperref}
% The repository URL is long enough to overflow a revtex column.  url.sty only
% breaks at a few characters; adding letters and digits to \UrlBreaks lets it
% break anywhere.  Done with url.sty itself rather than xurl, which is not in
% every TeX installation a journal might use.
\makeatletter
\g@addto@macro\UrlBreaks{%
  \do\a\do\b\do\c\do\d\do\e\do\f\do\g\do\h\do\i\do\j\do\k\do\l\do\m%
  \do\n\do\o\do\p\do\q\do\r\do\s\do\t\do\u\do\v\do\w\do\x\do\y\do\z%
  \do\A\do\B\do\C\do\D\do\E\do\F\do\G\do\H\do\I\do\J\do\K\do\L\do\M%
  \do\N\do\O\do\P\do\Q\do\R\do\S\do\T\do\U\do\V\do\W\do\X\do\Y\do\Z%
  \do\0\do\1\do\2\do\3\do\4\do\5\do\6\do\7\do\8\do\9\do\-\do\.}
\makeatother
\usepackage{placeins}
\hypersetup{hidelinks,
  pdftitle={A parity obstruction to degree-optimal syndrome extraction in bivariate bicycle codes},
  pdfauthor={Guangyi Li, Yu Gan, Zeguan Wu, Xueyue Zhang, Zheshen Zhang, Junyu Liu}}
\usepackage{booktabs}
\usepackage{enumitem}
\usepackage{flushend}

\renewcommand{\topfraction}{0.95}
\renewcommand{\bottomfraction}{0.9}
\renewcommand{\textfraction}{0.05}
\renewcommand{\floatpagefraction}{0.8}
\renewcommand{\dbltopfraction}{0.95}
\renewcommand{\dblfloatpagefraction}{0.8}
\setcounter{topnumber}{3}
\setcounter{dbltopnumber}{2}
\setcounter{totalnumber}{5}

\theoremstyle{plain}
\newtheorem{theorem}{Theorem}
\newtheorem{lemma}[theorem]{Lemma}
\newtheorem{proposition}[theorem]{Proposition}
\newtheorem{corollary}[theorem]{Corollary}
\newtheorem*{lemmaSI}{Lemma}
\newtheorem*{propSI}{Proposition}
\newtheorem*{remarkSI}{Remark}

\newcommand{\Zq}{\mathbb{Z}_q}
\newcommand{\F}{\mathbb{F}_2}
\newcommand{\col}{\mathsf{d}}          % collision element (kept distinct from the code distance d)
\newcommand{\colbar}{\bar{\mathsf{d}}} % its image under a twist
\newcommand{\wA}{w_A}                  % polynomial weights |A|, |B| (k and \ell are reserved for [[n,k,d]] and Z_\ell x Z_m)
\newcommand{\wB}{w_B}

\begin{document}

\title{A parity obstruction to degree-optimal syndrome extraction in bivariate bicycle codes}

\author{Guangyi Li}
\email{guangyi@umd.edu}
\affiliation{Joint Center for Quantum Information and Computer Science, University of Maryland, College Park, MD 20742, USA}
\affiliation{Department of Computer Science, University of Maryland, College Park, MD 20742, USA}
\author{Yu Gan}
\affiliation{Department of Computer Science, University of Pittsburgh, Pittsburgh, PA 15260, USA}
\author{Zeguan Wu}
\affiliation{Department of Computer Science, University of Pittsburgh, Pittsburgh, PA 15260, USA}
\author{Xueyue Zhang}
\affiliation{Department of Applied Physics and Applied Mathematics, Columbia University, New York, NY 10027, USA}
\author{Zheshen Zhang}
\affiliation{Department of Electrical Engineering and Computer Science, University of Michigan, Ann Arbor, MI 48109, USA}
\author{Junyu Liu}
\email[Corresponding author: ]{junyuliu@pitt.edu}
\affiliation{Department of Computer Science, University of Pittsburgh, Pittsburgh, PA 15260, USA}

\begin{abstract}
Syndrome extraction for a CSS code has depth at least the Tanner-graph
degree $\Delta$, attained by edge colouring if gate order is ignored;
correctness adds a parity constraint. We prove three
general obstructions to depth-$\Delta$
extraction and show that, for every two-block group-algebra code with
weight-three polynomials whose checks overlap in at most two qubits, no
correct single-round circuit with one ancilla per check and one CNOT per
Tanner edge has depth six: the minimum is seven, so the seven-layer cycle
of Bravyi \emph{et al.} is optimal for their seven bivariate bicycle
codes. This is the first analytic, family-wide impossibility result;
previous evidence was solver output. The hypothesis cannot be dropped: a
single orbit of four-qubit overlaps with a collision element of even order
and $2$-height at most eight makes depth six attainable, and a
$[[72,8,6]]$ code of this kind has a verified six-layer circuit. The
refutations carry DRAT certificates.
\end{abstract}

\maketitle

\section{Introduction}
\label{sec:intro}

Quantum low-density parity-check codes promise fault-tolerant memories
with constant overhead~\cite{BreuckmannEberhardt2021,Panteleev2022,
LeverrierZemor2022,Dinur2023}, and the bivariate bicycle (BB) codes of
Bravyi \emph{et al.}~\cite{Bravyi2024} are among the families closest to
implementation, with architectures~\cite{Yoder2025,Xu2024} and
variants~\cite{Voss2024,WangMueller2024,Berthusen2025,LiangChen2026} built
on them. A memory measures its stabilisers once per cycle, and the depth of
the syndrome-extraction circuit sets how many idle locations accumulate
errors between measurements. The published cycle uses seven layers of CNOT
gates for every member of the family~\cite{Bravyi2024}. Each data qubit
and each check has degree six in the Tanner graph, so six layers is a lower
bound for any circuit with one CNOT per Tanner edge. One CNOT layer is one
seventh of the published cycle and one idle location per qubit per cycle.
Throughout, optimality refers to the CNOT depth of a single interleaved
round with one ancilla per check and one CNOT per Tanner edge, without
extra ancillas, staggering across rounds, or connectivity constraints; it
makes no claim about circuit distance or logical error rate. This paper is
about the gap between six and seven.

Ignoring gate order, the bound is attainable: scheduling CNOTs is an
edge-colouring problem, K\H{o}nig's theorem colours a bipartite graph with
$\Delta$ colours~\cite{Konig1916}, and this is how depth-$\Delta$
extraction circuits are obtained for CSS codes~\cite{TDB2022};
depth-optimal synthesis of linear reversible circuits is studied in
Refs.~\cite{PMH2008,Jiang2020,MaslovZindorf2022,Kutin2007,Bataille2022}.
But a CNOT from an $X$-ancilla into a data qubit and a CNOT from that qubit
into a $Z$-ancilla do not commute: each $X$-first ordering on a shared
qubit adds a cross term, and a depth-$\Delta$ colouring is a correct
measurement only if the cross terms cancel. That order matters has been
known since the surface code~\cite{DKLP2002,Fowler2012,TomitaSvore2014};
for general CSS codes, Conrad \emph{et al.} observed that two CNOT pairs on
two shared qubits must agree in order~\cite{Conrad2018}, Zhang
\emph{et al.} stated the general condition, that the number of $X$-first
orderings over the shared qubits of any overlapping pair must be even, and
built a constraint solver around it~\cite{Zhang2026ASC}, as have other
automated schedulers~\cite{Liu2026AlphaSyndrome,Viszlai2026PropHunt}, and
Nguyen \emph{et al.} constructed schedules of depth $\Delta$ when one
polynomial weight is even and $\Delta+1$ when both are odd, leaving the
optimality of the latter to numerical evidence~\cite{Nguyen2026}. Bravyi
\emph{et al.} asked whether the depth of their cycle can be reduced and
reported numerical indications that the $[[72,12,6]]$ code admits no
depth-$6$ circuit~\cite[Sec.~10]{Bravyi2023arXiv}. Zhang \emph{et al.}
report that their CP-SAT compiler returned no depth-$6$ schedule for the
seven instances within a two-hour budget and describe this as first concrete
evidence; the runs were not proof-logged~\cite{Zhang2026ASC}. Depth-six
measurement circuits related to BB codes exist in other models: morphing
circuits change the code between measurements~\cite{ShawTerhal2025,
ShawTerhal2026}, and flag schemes change the gate
set~\cite{ChamberlandBeverland2018,ChaoReichardt2018}; neither belongs to
the published cycle's model.

We settle the question in that model: one ancilla per check, one CNOT per
Tanner edge, a single round, correctness meaning that the circuit
implements the ideal measurement. Our contributions are the following.
\begin{itemize}[leftmargin=*, itemsep=1pt, topsep=2pt]
\item[(i)] Three general obstructions to depth-$\Delta$ interleaved
extraction: a mod-4 condition on the degree profile of a CSS code
(Theorem~\ref{thm:mod4}, with its two-block and regular-CSS forms
Corollaries~\ref{cor:twoblock} and~\ref{cor:regular}), the impossibility of
a common layer set for two-block codes with both polynomial weights odd
(Theorem~\ref{thm:nomixed}), and a parity
condition on schedules twisted by a group homomorphism
(Theorem~\ref{thm:twisted}).
\item[(ii)] For two-block codes over any finite abelian group with
$|A|=|B|=3$ and pairwise distinct sums $a_s+b_t$---all seven BB instances
of Ref.~\cite{Bravyi2024} among them---the exact minimum depth of a correct
single-round interleaved circuit is $7$ (Theorem~\ref{thm:main},
Proposition~\ref{prop:seven}, Corollary~\ref{cor:bb}). The proof forces
every data qubit into a rigid pairing of adjacent layers and then counts
$X$-first corners, $(\Delta/2)^2$ of them, an odd number.
\item[(iii)] For codes with one collision orbit whose collision element
$\col$ has even order and $2$-height at most $8$, i.e.\ $\col\notin512G$,
the exact minimum depth is $6$ (Theorem~\ref{thm:six}), by a
group-independent template, a lifting lemma, and a reduction of the
template keys to a one-parameter family with a universal half-period
symmetry. The overlap hypothesis of (ii) therefore cannot be dropped.
\item[(iv)] Verified DRAT certificates for the seven codes of
Ref.~\cite{Bravyi2024} and eight further instances, an $\F$-simulation
check of every schedule we report, and a computational map of the
remaining regime.
\item[(v)] Depth-six codes with competitive parameters: a search over all
two-block codes of length up to $144$ yields, among others, $[[72,8,6]]$,
$[[108,12,6]]$, $[[144,16,6]]$ and $[[144,8,\le10]]$ codes with verified
six-layer circuits (Table~\ref{tab:depth6codes}); up to $n=48$ they match
the best depth-seven two-block codes, and from $n=72$ on they trail them by
one step in $k$ or in $d$.
\end{itemize}
We do not claim a criterion for general odd weights, a complete
characterisation of the coinciding-sum case, or any statement about
circuit-level distance. Nor do we claim optimality in amortised or
staggered models: the published cycle takes $8C+1$ time steps for $C$
rounds~\cite{Yoder2025}, and left--right circuits reach amortised depth
$\Delta$ by overlapping rounds~\cite{Strikis2026}. Lower bounds for
circuits with CNOTs outside the Tanner graph or with extra ancillas are
open.

\section{Model and correctness condition}
\label{sec:model}

\paragraph*{The cross term.}
Let $H_X,H_Z$ be the check matrices of a CSS code, $\mathbf{x}$ the data
register, $\mathbf{y}$ the $X$-ancillas (prepared in $|+\rangle$) and
$\mathbf{z}$ the $Z$-ancillas (prepared in $|0\rangle$). Each Tanner edge
is executed by one CNOT, controlled by the $X$-ancilla for $X$-edges and by
the data qubit for $Z$-edges; there are no other gates and no further
ancillas, and the ancillas are measured after the last layer. A
\emph{schedule} $\pi$ assigns layers $1,\dots,D$ to the edges so that no
qubit is in two CNOTs of one layer; $D$ is its depth. The circuit implements
\begin{equation}
(\mathbf{x},\mathbf{y},\mathbf{z})\mapsto
\bigl(\mathbf{x}+H_X^{\mathsf T}\mathbf{y},\;\mathbf{y},\;
\mathbf{z}+H_Z\mathbf{x}+M_\pi\mathbf{y}\bigr),
\label{eq:se-map}
\end{equation}
where $(M_\pi)_{ji}$ is the number, modulo $2$, of shared qubits of $X_i$
and $Z_j$ on which the $X$-edge precedes the $Z$-edge.

Correctness means that for every input state of the data register the
circuit, followed by the $X$-basis measurement of $\mathbf{y}$ and the
$Z$-basis measurement of $\mathbf{z}$, realises the ideal joint projective
measurement of the commuting checks $\{X_i\}\cup\{Z_j\}$: the recorded
outcome associated with each ancilla is the eigenvalue of its designated
check, and the post-measurement data state is the projected state.
Redundancy-dependent classical recombination of the labelled outcomes is
not included in this definition.

Propagating the measured Pauli operators backwards through the CNOTs, and
using the even overlap of each commuting $X$/$Z$ check pair in the second
identity (Supplementary Section~S2), gives
\begin{equation}
\begin{aligned}
U^\dagger Z_{z_j} U
  &= Z_{z_j}\,\overline{Z}_j
     \prod_i Z_{a_i}^{(M_\pi)_{ji}}, \\
U^\dagger X_{a_i} U
  &= X_{a_i}\,\overline{X}_i
     \prod_j X_{z_j}^{(M_\pi)_{ji}} .
\end{aligned}
\label{eq:heisenberg}
\end{equation}
Here $\overline{Z}_j,\overline{X}_i$ are the checks on the data register
and $a_i,z_j$ the ancillas. If $M_\pi=0$, these are the ideal measurement
operators, since each ancilla is a $+1$ eigenstate of its own measured
Pauli.

Conversely, suppose that $M_\pi\neq0$, and let $u_j$ be a nonzero row of
$M_\pi$. Write
\[
Z_a^{u_j}:=\prod_i Z_{a_i}^{(u_j)_i}.
\]
For any input data state $\rho$, the marginal expectation of the recorded
outcome $s_j\in\{\pm1\}$ of $Z_{z_j}$ is
\[
\begin{aligned}
\mathbb{E}_{\rho}[s_j]
&=
\operatorname{Tr}\!\left[
  \bigl(\rho\otimes |+\rangle\!\langle+|_a
  \otimes |0\rangle\!\langle0|_z\bigr)
  U^\dagger Z_{z_j}U
\right] \\
&=
\operatorname{Tr}(\rho\,\overline{Z}_j)\,
{}_z\langle0|Z_{z_j}|0\rangle_z\,
{}_a\langle+|Z_a^{u_j}|+\rangle_a
=0,
\end{aligned}
\]
because $u_j\neq0$ and hence
${}_a\langle+|Z_a^{u_j}|+\rangle_a=0$. Since $s_j$ is binary, its marginal
distribution is therefore uniform. In particular, on a simultaneous
eigenstate of the commuting checks, the ideal outcome of
$\overline{Z}_j$ is deterministic whereas the recorded outcome $s_j$ is
uniform. Thus the circuit realises the stipulated ideal measurement if
and only if $M_\pi=0$.

The matrix $M_\pi$ is therefore the ancilla-to-ancilla contamination
matrix, and a schedule with $M_\pi=0$ is called \emph{proper}. The maximum
degree $\Delta$ of the Tanner graph bounds $D$ from below, and
K\H{o}nig's theorem gives a schedule of depth exactly $\Delta$ that need
not be proper (Fig.~\ref{fig:setup}; Supplementary Section~S1 records
this as a compiler with an exact depth certificate for commuting gate
sets).
\paragraph*{Two-block codes.}
Let $G$ be a finite abelian group of order $N$, written additively, and
$A=\{a_1,\dots,a_{\wA}\}$, $B=\{b_1,\dots,b_{\wB}\}\subset G$. The two-block
code $\mathrm{QC}(A,B)$~\cite{KovalevPryadko2013,PanteleevKalachev2021,
LinPryadko2024,WangLinPryadko2023} has data qubits $L_g,R_g$ and checks
$X_g,Z_g$, $g\in G$, with
\begin{equation}
X_g:\ L_{g+a_s},\,R_{g+b_t};\qquad Z_g:\ L_{g-b_t},\,R_{g-a_s};
\label{eq:twoblock}
\end{equation}
so $H_X=[A\,|\,B]$, $H_Z=[B^{\mathsf T}\,|\,A^{\mathsf T}]$ in the
notation of Ref.~\cite{Bravyi2024}, the BB codes being the case
$G=\mathbb{Z}_\ell\times\mathbb{Z}_m$ with $\wA=\wB=3$. The \emph{weights}
$\wA=|A|$ and $\wB=|B|$ are the numbers of monomials of the two
polynomials; every check has weight $\wA+\wB$, every $L$-qubit has $\wA$
$X$- and $\wB$ $Z$-neighbours, every $R$-qubit the reverse, and
$\Delta=\wA+\wB$. Unless stated otherwise $\wA=\wB=3$ and $\Delta=6$. We
write $H\le G$ for the subgroup generated by the differences $a_s-a_{s'}$
and $b_t-b_{t'}$; the connected components of the Tanner graph are the
cosets of $H$ (Corollary~\ref{cor:twoblock}). An edge carries the
\emph{label} $s$ or $t$ of the element $a_s$ or $b_t$ that defines it. $X_g$ and $Z_{g'}$ share $L_{g+a_s}$ and $R_{g+b_t}$ exactly
when $g'-g=a_s+b_t$; with $S(\delta)=\#\{(s,t):a_s+b_t=\delta\}$ they
overlap in $2S(g'-g)$ qubits. The code has \emph{distinct sums} (class D)
if $S\le1$ everywhere, so that every overlap has size $0$ or $2$; in the
polynomial notation this is the condition that the nine products $a_sb_t$
are distinct, and it holds for all seven codes of Ref.~\cite{Bravyi2024}.
It is of \emph{class C1} if $S(\delta_0)=2$ for exactly one $\delta_0$ and
$S\le1$ elsewhere; the two label pairs $(s,t),(s',t')$ with sum $\delta_0$
form the \emph{collision}, with \emph{collision element}
$\col=a_{s'}-a_s=b_t-b_{t'}$, and the $N$ pairs $(X_g,Z_{g+\delta_0})$
overlap in four qubits. Here $2^jG=\{2^jg:g\in G\}$, $\operatorname{ord}(\col)$
is the order of $\col$ in $G$, and for $\col$ of even order the
\emph{$2$-height} $h(\col)$ is the largest $h$ with $\col\in2^hG$; it is
finite exactly when the order is even, because the elements lying in every
$2^jG$ form the odd-order subgroup of $G$. Class C2 has two collisions or
one triple coincidence (Supplementary Section~S6).

\paragraph*{Corners and counting.}
A \emph{corner} is a data qubit $q$ with an $X$-edge $e$ and a $Z$-edge $e'$
at $q$; its orientation is $1$ if $\pi(e)<\pi(e')$ and $0$ otherwise. The
entry of $M_\pi$ for a pair is the orientation sum over its shared corners,
so $M_\pi=0$ reads
\begin{equation}
\sum_{q\in\operatorname{supp}X_g\cap\operatorname{supp}Z_{g'}}
o\bigl(q,e_g(q),e'_{g'}(q)\bigr)\equiv0\pmod 2
\label{eq:parity}
\end{equation}
for every overlapping pair, the constraint of
Refs.~\cite{Conrad2018,Zhang2026ASC,Nguyen2026}. Under distinct sums each
pair has one $L$-corner and one $R$-corner, which must have equal
orientation; this pairs the corners of the two blocks by a
label-preserving bijection $\Phi$ (Supplementary Section~S2). A four-qubit
overlap allows two $X$-first and two $Z$-first corners. For an $X$-check
$g$ let $P_g,V_g$ be the layer sets of its $L$- and $R$-edges; at depth
$\Delta$, $P_g\sqcup V_g=[\Delta]:=\{1,\dots,\Delta\}$ and every qubit and
check has one edge per layer. Let $\gamma_q(i)$ be the number of $Z$-edges
at $q$ in layers above $i$, and
\begin{equation}
F_L(g)=\sum_{i\in P_g}\gamma(i),\qquad F_R(g)=\sum_{i\in V_g}\gamma(i),
\label{eq:FLFR}
\end{equation}
evaluated at the qubit carrying each edge: the $X$-first corners of $X_g$ on
the two blocks. Each such corner lies in one overlapping pair, so
Eq.~\eqref{eq:parity} gives $F_L+F_R\equiv0\pmod 2$ for any overlap
pattern, and $F_L=F_R$ under distinct sums. A schedule is
\emph{translation-invariant} if $\pi$ depends only on label and edge type,
\emph{twisted} by a homomorphism $\varepsilon:G\to\Zq$ if it depends only
on label, type and $\varepsilon(g)$, and a layer is \emph{mixed} if it
holds $L$-edges of some $X$-checks and $R$-edges of others.

\paragraph*{Terminology.}
\emph{Layer $i$ is above layer $j$} means $i>j$; ``$X$-first'' means the
$X$-edge is in the lower layer. A \emph{layering} is any assignment of
layers with no qubit in two CNOTs of one layer; a \emph{schedule} is a
layering of depth $D$, and it is \emph{proper} if $M_\pi=0$.
\emph{Interleaved} means that $X$- and $Z$-ancillas share layers;
\emph{single-round} means one measurement of every check. An \emph{orbit}
of overlapping pairs is the set $\{(X_g,Z_{g+\delta}):g\in G\}$ for fixed
$\delta$. A schedule has \emph{period} $q$ if it is twisted by a
homomorphism onto $\Zq$ with $q>1$ dividing $N$; a found schedule has
\emph{no proper period} if it is not twisted by any such homomorphism, a
finite check performed on that schedule, which by itself says nothing
about other schedules of the same code. $\mathrm{QC}$ stands for
quasi-cyclic. The letters $k$ and $d$ are reserved for the code parameters
$[[n,k,d]]$ and $\ell,m$ for the BB group $\mathbb{Z}_\ell\times\mathbb{Z}_m$;
polynomial weights are $\wA,\wB$ and the collision element is $\col$.

\begin{figure*}[t]
\centering
\includegraphics[width=\textwidth]{fig1_degree_bound_to_cross_term.pdf}
\caption{\textbf{Degree bound versus correctness.} Left: a patch of the
$[[72,12,6]]$ Tanner graph drawn as two plates, checks above ($X$ blue,
$Z$ orange) and data qubits below; one $X$-check $X_g$ and one $Z$-check
$Z_{g'}$ sharing an $L$-qubit and an $R$-qubit are highlighted with their
twelve edges, the four on the shared qubits numbered by their layer in the
depth-$7$ schedule of Proposition~\ref{prop:seven} (edge colour = layer).
Right: the same four CNOTs on the four wires of the pair; faded gates lead
to the checks' other qubits. Top: in the published cycle the $X$-edge
precedes the $Z$-edge on both shared qubits, so the $X$-ancilla value $y$
enters the $Z$-ancilla twice and cancels ($o_L=o_R=1$). Bottom: a proper
$6$-edge-colouring of the full $[[72,12,6]]$ Tanner graph, restricted to
the pair, in which $y$ enters once, leaving $M_\pi\neq0$.}
\label{fig:setup}
\end{figure*}

\section{Results}
\label{sec:results}

Fig.~\ref{fig:mechanism} shows the mechanism; the full dependency map is
Supplementary Fig.~\ref{fig:fullmap}, and proofs not given in Methods are
in Supplementary Sections~S3--S5.

\subsection{Three general obstructions}

\begin{theorem}[mod-4 obstruction]
\label{thm:mod4}
Let every data qubit $v$ of a CSS code have $r_v$ $X$-neighbours and $s_v$
$Z$-neighbours with $r_v+s_v=\Delta$, and let every $X$-check have weight
$\Delta$. If $\sum_v r_vs_v\not\equiv0\pmod 4$ over some connected
component of the Tanner graph, no proper schedule of depth $\Delta$
exists.
\end{theorem}

Any depth-$\Delta$ layering of the component has exactly
$\frac12\sum_vr_vs_v$ $X$-first corners: explicitly,
$\sum_vN_v=\tfrac{\Delta-1}{2}\sum_vr_v-\sum_v\binom{r_v}{2}
=\tfrac12\sum_vr_vs_v$ (Methods), and properness makes every pair's count
even. Only the weight of the $X$-checks and the degrees of the data qubits
enter; the $Z$-check weights are unconstrained. For a two-block code every
data qubit has $r_vs_v=\wA\wB$ and the component structure is explicit.

\begin{corollary}[two-block form]
\label{cor:twoblock}
Let $H\le G$ be the subgroup generated by the differences $a_s-a_{s'}$ and
$b_t-b_{t'}$. The component of the Tanner graph containing $X_g$ consists of
the checks $X_{g+H}$, $Z_{g+a_1+b_1+H}$ and the data qubits $L_{g+a_1+H}$,
$R_{g+b_1+H}$, so it has $|H|$ checks of each type and $2|H|$ data qubits.
If $|H|\,\wA\wB$ is odd, $\mathrm{QC}(A,B)$ has no proper schedule of depth
$\wA+\wB$.
\end{corollary}

The proof is in Methods. For a connected Tanner graph ($H=G$) the condition
is that $N$, $\wA$ and $\wB$ are all odd; for a BB code with a connected
Tanner graph and $2\ell m$ data qubits, $\sum_vr_vs_v=18\ell m$ and the
condition reads $\ell m$ odd. The Tanner graphs of the seven codes of
Ref.~\cite{Bravyi2024} are connected, and only $[[90,8,10]]$, with
$\ell m=45$, is excluded this way. The $\mathbb{Z}_6$ instance of
Section~\ref{sec:sharpness}, $A=B=\{0,2,4\}$, has $H=\{0,2,4\}$ and is
excluded although $|G|$ is even, through components of six data qubits
with $\sum_vr_vs_v=54$. The same subgroup $H$ reappears in
Theorem~\ref{thm:twisted}, whose $M$ is $\varepsilon(H)$. Beyond two-block
codes the regular case reads as follows.

\begin{corollary}[regular CSS codes]
\label{cor:regular}
Let every data qubit of a CSS code have exactly $w$ $X$- and $w$
$Z$-neighbours and every $X$-check have weight $2w$. If a connected
component of the Tanner graph contains $n_0$ data qubits with $w$ odd and
$n_0\equiv2\pmod 4$, no proper schedule of depth $2w$ exists.
\end{corollary}

Indeed $\sum_vr_vs_v=n_0w^2\equiv2\pmod4$. The hypergraph product of a
$(3,3)$-regular classical code of odd length $n_c$ with itself has
$2n_c^2$ data qubits of type $(3,3)$ and weight-$6$ checks; for the
circulant code $1+x+x^3$ of length $7$ it is a connected $[[98,18]]$ code,
for length $15$ a connected $[[450,32]]$ code, and depth $6$ is impossible
for both.

\begin{theorem}[no common layer set]
\label{thm:nomixed}
For a two-block code with $|A|=\wA$, $|B|=\wB$ and $\wA\wB$ odd, every
proper schedule of depth $\wA+\wB$ contains a mixed layer; in particular no
translation-invariant proper schedule of depth $\Delta$ exists.
\end{theorem}

With a common $L$-layer set $P$ every $L$-qubit has $X$-layers $P$ and
$Z$-layers $[\Delta]\setminus P$, every $R$-qubit the reverse, and
$F_L+F_R=|P|\,|[\Delta]\setminus P|=\wA\wB$ must be even. This is why the
translation-invariant construction of Ref.~\cite{Nguyen2026} reaches
$\Delta$ only for one even weight.

\begin{theorem}[twisted schedules]
\label{thm:twisted}
Let $\mathrm{QC}(A,B)$ be a two-block code with an arbitrary overlap
pattern, $|A|=\wA$, $|B|=\wB$, let $\varepsilon:G\to\Zq$ be a
homomorphism, and let $M=\varepsilon(H)\le\Zq$ be generated by the images
of the differences $a_s-a_{s'}$, $b_t-b_{t'}$. If a proper schedule of depth
$\wA+\wB$ twisted by $\varepsilon$ exists, then $|M|\,\wA\wB$ is even; for
$\wA=\wB=3$, $|M|$ is even and no such schedule exists for odd $q$.
\end{theorem}

The classes in one coset of $M$ form a closed subsystem in which
$\sum_c(F_L(c)+F_R(c))=|M|\wA\wB$ while every summand is even. None of the
three theorems excludes a proper schedule that breaks all symmetry; that
is the main theorem.

\subsection{Exact depth seven}
\label{sec:main}

For a proper depth-$\Delta$ schedule let $a$ and $b$ be the lowest and
highest layers in which any $X$-check has an $L$-edge.

\begin{lemma}[support symmetry]
\label{lem:support}
For a two-block code with distinct sums and any $|A|,|B|$, every proper
schedule of depth $\Delta$ has $a+b=\Delta+1$, and an $X$-check with an
$L$-edge in layer $a$ has one in layer $b$, and conversely.
\end{lemma}

From here on let $|A|=|B|=3$ with distinct sums. Lemma~\ref{lem:support}
leaves $(a,b)\in\{(3,4),(2,5),(1,6)\}$: the first is empty, since an
$X$-check needs three $L$-layers in $[a,b]$; the second is excluded by
counting the corners of the layer-$2$ and layer-$5$ edges (Supplementary
Theorem~\ref{thm:S3}). The third is settled by two lemmas.

\begin{lemma}[rigidity]
\label{lem:rigidity}
For a code with $|A|=|B|=3$ and distinct sums, in every proper schedule of
depth $6$ with $(a,b)=(1,6)$, every data qubit carries exactly one $X$-edge
and one $Z$-edge in each of the layer pairs $\{1,2\}$, $\{3,4\}$,
$\{5,6\}$.
\end{lemma}

\begin{lemma}[parity]
\label{lem:parity}
Let $|A|=|B|=w$, $\Delta=2w$. If a depth-$\Delta$ schedule has the pair
structure of Lemma~\ref{lem:rigidity} at every data qubit, then
$F_L(g)+F_R(g)=w^2$ for every $X$-check $g$.
\end{lemma}

Under the pair structure the number of $Z$-edges above an $X$-edge depends
only on its layer, $\gamma(i)=w-\lceil i/2\rceil+\mathbb{1}[i\text{ odd}]$,
and $\sum_{i\le2w}\gamma(i)=w^2$. The rigidity proof starts from an
$X$-check with $L$-edges in layers $1$ and $6$, follows $\Phi$ through its
$R$-edges, and shows layer by layer that any other arrangement produces a
corner whose partner cannot be placed (Supplementary Section~S4).

\begin{theorem}[main theorem]
\label{thm:main}
Let $G$ be a finite abelian group and $A,B\subset G$ three-element sets
with the nine sums $a_s+b_t$ pairwise distinct. Then $\mathrm{QC}(A,B)$
admits no proper single-round interleaved schedule of CNOT depth $6$.
\end{theorem}

\begin{proof}
Only $(a,b)=(1,6)$ remains. Lemma~\ref{lem:rigidity} gives the pair
structure, so Lemma~\ref{lem:parity} gives $F_L(g)+F_R(g)=9$ for every
$X$-check. Distinct sums give $F_L(g)=F_R(g)$, hence
$F_L(g)+F_R(g)=2F_L(g)\equiv0\pmod 2$, contradicting $9$.
\end{proof}

\begin{proposition}[uniform upper bound]
\label{prop:seven}
For every two-block code with $|A|=|B|=3$, the translation-invariant
schedule of Table~\ref{tab:depth7main} is a proper schedule of depth $7$.
Hence the minimum depth of every such code lies in $\{6,7\}$.
\end{proposition}

\begin{proof}
The $L$-corner of $X_g$ with labels $(s,t)$ has orientation
$\mathbb{1}[\pi_{XL}(s)<\pi_{ZL}(t)]$ and the $R$-corner
$\mathbb{1}[\pi_{XR}(t)<\pi_{ZR}(s)]$. Since $\pi_{XL}=(1,2,7)$ places
$s=1,2$ below and $s=3$ above every entry of $\pi_{ZL}=(5,3,4)$, and
$\pi_{ZR}=(6,7,1)$ against $\pi_{XR}=(4,3,5)$ does the same, both
orientations equal $\mathbb{1}[s\le2]$ for all nine label pairs. Every
overlapping pair consists of corners indexed by label pairs, so its
orientation sum is even and Eq.~\eqref{eq:parity} holds for any overlap
pattern.
\end{proof}

\begin{table}[t]
\caption{\textbf{The uniform depth-$7$ schedule of
Proposition~\ref{prop:seven}.} Layers by check type, block and label.}
\label{tab:depth7main}
\begin{ruledtabular}
\begin{tabular}{lccc}
 & label 1 & label 2 & label 3 \\
\colrule
$X$-check, $L$-edges & 1 & 2 & 7 \\
$X$-check, $R$-edges & 4 & 3 & 5 \\
$Z$-check, $L$-edges & 5 & 3 & 4 \\
$Z$-check, $R$-edges & 6 & 7 & 1 \\
\end{tabular}
\end{ruledtabular}
\end{table}

\begin{corollary}
\label{cor:bb}
For every bivariate bicycle code with pairwise distinct products $a_sb_t$,
including all seven codes of Ref.~\cite{Bravyi2024}, the minimum CNOT
depth of a proper single-round interleaved schedule is exactly $7$.
\end{corollary}

The lower bound is Theorem~\ref{thm:main} and the upper bound
Proposition~\ref{prop:seven}, whose schedule was also verified by $\F$
simulation on the seven codes of Ref.~\cite{Bravyi2024} and on $40$ random
instances (Supplementary Section~S7). The decomposition is more general
than the case it settles: whenever a schedule of a code with $\wA=\wB$ has
the pair structure, Lemma~\ref{lem:parity} and $F_L=F_R$ contradict each
other exactly when $\Delta\equiv2\pmod 4$; what does not yet extend to
$\wA=\wB\ge5$ is the rigidity step. Theorems~\ref{thm:S1} and~\ref{thm:S2}
give depth $\Delta$ unattainable when one weight is $1$ and the other odd,
and when $\wA,\wB$ are odd and $a=(\wB+1)/2$.

\begin{figure*}[t]
\centering
\includegraphics[width=\textwidth]{fig2_why_six_layers_fail.pdf}
\caption{\textbf{Why six layers fail, and when they do not.} Left
staircase, distinct sums: from $M_\pi=0$, Lemma~\ref{lem:support} confines
the $X$-edges on $L$ to $[a,b]$ with $a+b=7$ and excludes the cases $(3,4)$
and $(2,5)$; Lemma~\ref{lem:rigidity} forces one $X$- and one $Z$-edge into
each adjacent layer pair; Lemma~\ref{lem:parity} then gives $F_L+F_R=9$,
odd, while $F_L=F_R$ requires it even (Theorem~\ref{thm:main}), so depth $7$
is optimal (Proposition~\ref{prop:seven}). Right staircase, one repeated
sum: four shared corners can cancel, and a class schedule on
$\mathbb{Z}_2$ lifted along $\varepsilon$ (Lemma~\ref{lem:lift},
Proposition~\ref{prop:template}) attains depth $6$ whenever the collision
element has even order and $2$-height at most $8$ (Theorem~\ref{thm:six}). The ruler marks the minimum depth reached by each
staircase; cell colours in the class strip are the layer scale of
Fig.~\ref{fig:setup}.}
\label{fig:mechanism}
\end{figure*}

\subsection{Exact depth six}
\label{sec:six}

Fix $q$, a key $\kappa=(\alpha_1,\alpha_2,\alpha_3;\beta_1,\beta_2,\beta_3)
\in\Zq^6$ and a designated pair $\{(s,t),(s',t')\}$ with
$\alpha_s+\beta_t=\alpha_{s'}+\beta_{t'}$; write $\colbar=\alpha_{s'}-\alpha_s$.
A \emph{class schedule} assigns to every class $c\in\Zq$, label and edge
type a layer in $[6]$ such that (i)~the six edges incident to the
$X$-check of class $c$ occupy six distinct layers, and so do the six edges
incident to the $Z$-check of class $c$; (ii)~at the $L$-qubit of class
$c$, the $X$-edges of label $s$ from class $c-\alpha_s$ and the $Z$-edges
of label $t$ from class $c+\beta_t$ occupy distinct layers, and likewise
on the $R$ block; (iii)~for every class and every label pair other than
the designated two, the $L$-corner of $X_c$ with labels $(s,t)$ and the
$R$-corner with the same labels, at the $Z$-check of class
$c+\alpha_s+\beta_t$, have equal orientation; (iv)~for the designated pairs
the four orientations have even sum. The constraints are indexed by
labels, not by which projected sums coincide in $\Zq$.

\begin{lemma}[lifting]
\label{lem:lift}
Let $\mathrm{QC}(A,B)$ be of class C1 with collision $\{(s,t),(s',t')\}$,
$\varepsilon:G\to\Zq$ a homomorphism, and $T$ a class schedule for the key
$(\varepsilon(a_s);\varepsilon(b_t))$ and that collision. Then
$\pi(e)=T(\varepsilon(g),\mathrm{label}(e),\mathrm{type}(e))$ is a proper
schedule of depth $6$ of $\mathrm{QC}(A,B)$.
\end{lemma}

Every constraint of $\mathrm{QC}(A,B)$ is indexed by a labelled object: a
check (i), a data qubit (ii), an overlapping pair together with its unique
label pair (iii), or the collision pair together with its two label pairs
(iv). The edges at $L_h$ come from $X_{h-a_s}$ and $Z_{h+b_t}$, of classes
$\varepsilon(h)-\alpha_s$ and $\varepsilon(h)+\beta_t$, so the conflict
constraints of the code are the class constraints (i)--(ii) at
$\varepsilon(g)$, $\varepsilon(h)$, $\varepsilon(r)$ evaluated on
$T(\varepsilon(\cdot))$; the overlapping pairs of a class-C1 code are
exactly the labelled pairs, the two collision pairs sharing one $Z$-check,
so Eq.~\eqref{eq:parity} for the pair $(X_g,Z_{g+a_s+b_t})$ is the class
constraint (iii) for $(\varepsilon(g),s,t)$, or (iv) for the collision.
Coincidences among projected sums $\alpha_s+\beta_t$ create no constraint
in $\mathrm{QC}(A,B)$, whose pairs are determined by the sums in $G$, and
remove none, because (iii) is imposed per class and label pair rather than
per class-level pair.

\begin{proposition}[templates, computer-verified]
\label{prop:template}
For $q\in\{2,4,8\}$ every key with $\colbar\neq0$ admits a class schedule,
and for $q=2^j$ with $4\le j\le9$ every key with $\colbar=q/2$ admits one.
For every member of the family~\eqref{eq:family} and every $j\le9$ a class
schedule exists that satisfies $T(c+q/2,e)=\sigma_{\mathrm{blk}(e)}(T(c,e))$
with the fixed involutions $\sigma_L=(13)(24)(56)$ on $L$-edges and
$\sigma_R=(12)(36)(45)$ on $R$-edges.
\end{proposition}

For $q\in\{2,4,8\}$ this was established by exhaustive enumeration: every
normalised key and designated collision was decided by SAT, one class
schedule stored per admissible key, and every stored schedule checked
against (i)--(iv) and by $\F$ simulation of its pullback to a realising
class-C1 code (Supplementary Section~S5, which also records the $q=2$
template and its pullback to a $[[28,6,4]]$ code over $\mathbb{Z}_{14}$).
For $q\ge16$ the keys with $\colbar=q/2$ reduce, by the symmetries of the
class system, to the one-parameter family~\eqref{eq:family} below
(Methods); every member of the family was decided satisfiable by SAT for
$q\le512$, without any assumption for $q\le32$ and under the half-period
symmetry of the statement for $q\ge64$, where the symmetry turns instances
that take up to half an hour into instances decided in a second. Odd $q$
is excluded analytically by
Theorem~\ref{thm:twisted}; $q=3$ and $q=5$ were also enumerated and admit
nothing, and $q=6$ was enumerated to test the refinement discussed in
Section~\ref{sec:sharpness}. In the enumeration, no key with $\colbar=0$
admits a class schedule.

The moduli are chosen by the $2$-height of the collision element.

\begin{lemma}[$2$-height]
\label{lem:height}
If $\col\in G$ has even order and $2$-height $h$, there is a homomorphism
$\varepsilon:G\to\mathbb{Z}_{2^{h+1}}$ with $\varepsilon(\col)=2^h$.
\end{lemma}

The proof (Supplementary Section~S5) splits $G/2^{h+1}G$ into cyclic
$2$-groups and picks a factor on which $\col$ survives.

\begin{theorem}[depth six]
\label{thm:six}
Let $\mathrm{QC}(A,B)$ be of class C1 with a collision element $\col$ of
even order and $2$-height $h(\col)\le8$, equivalently $\col\notin512G$.
Then the minimum CNOT depth of a proper single-round interleaved schedule
is exactly $6$.
\end{theorem}

\begin{proof}
Lemma~\ref{lem:height} gives $\varepsilon:G\to\Zq$ with $q=2^{h+1}\le512$
and $\varepsilon(\col)=q/2$, so the key $(\varepsilon(a_s);\varepsilon(b_t))$
has $\colbar=q/2$ and, after the normalisation of Methods, is one of the
keys of Proposition~\ref{prop:template}. Lemma~\ref{lem:lift} pulls the
class schedule back to a proper schedule of depth $6$; the degree bound
gives the lower bound.
\end{proof}

Every element of even order has a finite $2$-height, so the only gap
between Theorem~\ref{thm:six} and the conjecture that even order suffices
is the modulus bound. By the reduction in Methods, the keys with
$\colbar=q/2$ at $q=2^j$ reduce, up to the symmetries of the class system
and a decomposition into subsystems of smaller modulus, to the
one-parameter family
\begin{equation}
\kappa_q(1,b)=(0,\,q/2,\,1;\;0,\,q/2,\,b),\qquad b\in\Zq ,
\label{eq:family}
\end{equation}
and the conjecture is equivalent to the statement that every member of
this family admits a class schedule for every $j$; it holds for $j\le9$.
Every class schedule found for the family has the half-period symmetry of
Proposition~\ref{prop:template}, with the same two involutions for every
$j$ and $b$; the symmetry fixes the layers on the classes $c\ge q/2$ in
terms of those on $c<q/2$, and is the natural starting point for an
analytic template (Supplementary Section~S5). All $265$ class-C1 instances with $\operatorname{ord}(\col)$ even in our
data set have $h(\col)\le3$ and are covered by Theorem~\ref{thm:six}
(Supplementary Table~\ref{tab:twoheight}).

\subsection{The remaining regime}
\label{sec:sharpness}

Fig.~\ref{fig:sharpness} maps the depth-$6$ decision problem over $664$
instances with a recorded depth-$6$ outcome: the seven codes of
Ref.~\cite{Bravyi2024}, the $\mathbb{Z}_6$ instance of
Theorem~\ref{thm:mod4}, and $656$ sampled instances on cyclic groups of
order $10$--$24$, $28$, $30$, $32$ and $36$ (sampling and equivalence in
Supplementary Section~S6). For even $|G|\le8$ the nine sums cannot be
distinct, and in an exhaustive sweep of all $1441$ normalised instances
($33$ inequivalent) $1440$ admit a proper depth-$6$ schedule, the
exception being excluded by Theorem~\ref{thm:mod4}. Order $9$ is the
smallest order with distinct-sums instances: up to equivalence there are
three, one on $\mathbb{Z}_9$ and two on $\mathbb{Z}_3\times\mathbb{Z}_3$,
and all three are unsatisfiable, as Theorem~\ref{thm:main} requires; so is
every sampled distinct-sums instance of larger order. For $\wA=\wB=3$,
Theorem~\ref{thm:twisted} excludes every twisted schedule of odd period, so
a proper schedule is either twisted with even period or has no period.
When $\operatorname{ord}(\col)$ is odd, the found schedules are of both
kinds: of $333$ such class-C1 instances, $211$ are excluded by
Theorem~\ref{thm:mod4}; of the remaining $122$, $30$ admit depth $6$, $62$
are refuted by the solver, $6$ of them with verified DRAT certificates and
$56$ without proof logging, and $30$ were undetermined within the cap. Of
the $30$ satisfiable instances, $24$ have a schedule twisted with an even
period, consistent with Theorem~\ref{thm:twisted}: the relevant parity is
that of $|M|$, not of $\operatorname{ord}(\col)$; for the remaining $6$ no
twisted proper schedule exists for any $q\mid N$ (the search over all
proper divisors of $N$ completed without a cap), so every proper depth-$6$
schedule of those six codes is unperiodic (Supplementary
Fig.~\ref{fig:periods}). The refinement that a class schedule exists
exactly when $\operatorname{ord}(\colbar)$ is even is false, with $468$
admissible keys at $q=6$ having $\operatorname{ord}(\colbar)=3$. A
criterion for the odd-order regime is open; Section~\ref{sec:codes} shows
that it contains codes of interest.

\begin{figure*}[t]
\centering
\includegraphics[width=\textwidth]{fig3_where_depth6_is_attainable.pdf}
\caption{\textbf{Where depth $6$ is attainable.} \textbf{(a)}~The nine
sums $a_s+b_t$ of a two-block code; a repeated sum (rose link) defines the
collision element $\col$ and turns one orbit of two-qubit overlaps into
four-qubit overlaps. \textbf{(b)}~Minimum depth of a proper schedule by
collision structure. Distinct sums: $7$ for every code
(Theorem~\ref{thm:main}; $\blacklozenge$ nine instances also
DRAT-certified). One collision with $\col$ of even order: $6$ for every
code with $h(\col)\le8$ (Theorem~\ref{thm:six}), which covers all $265$
tested instances, the two with $\col\in8G$ through the $q=16$ templates
(they were also found as twisted witnesses). Odd order: of $122$
instances not excluded by Theorem~\ref{thm:mod4}, $30$ reach depth $6$,
$62$ are refuted by the solver and $30$ are open; $211$ further instances
are excluded by Theorem~\ref{thm:mod4}. Fill encodes settled status
(theorem or certificate), outline the solver status, hatching a capped
run. The $56$ solver-reported refutations carry no proof certificate and
are marked as computational indication; $6$ further instances in that cell
are DRAT-certified. The plotted classes D and C1 are drawn from the
$664$-instance dataset; the $28$ class-C2 and $3$ subgroup-type instances
are in Supplementary Table~\ref{tab:h2}.}
\label{fig:sharpness}
\end{figure*}

\subsection{Depth-six codes with competitive parameters}
\label{sec:codes}

Theorem~\ref{thm:six} says nothing about distance, so we searched. For
every abelian group of order $N\in\{18,24,27,36,45,54,72\}$ we classified
the two-block codes with $\wA=\wB=3$---all pairs $(A,B)$ with $0\in A$,
$0\in B$ for $N\le36$, $4\times10^5$ random pairs per group for
$N\ge45$---by the multiset of sums, computed $k$, and estimated $d$ for the
orbit representatives under $\mathrm{Aut}(G)$, translations and block swap
by information-set decoding, confirming the bound by an exhaustive search
for a lighter logical operator whenever that was feasible (Methods).
Table~\ref{tab:depth6codes} lists the Pareto-optimal $(k,d)$ pairs by
length and collision structure; Supplementary Section~S8 lists the codes.

\begin{table*}[t]
\caption{\textbf{Best parameters found for two-block codes with
$\wA=\wB=3$, by length and collision structure.} Pareto-optimal $(k,d)$
pairs with $d\ge4$ over all abelian groups of order $N=n/2$ (exhaustive
over all $(A,B)$ for $N\le36$; $4\times10^5$ sampled pairs per group for
$N\ge45$). Distinct sums: depth $7$ by Theorem~\ref{thm:main}. Class C1
with $\operatorname{ord}(\col)$ even: depth $6$ by Theorem~\ref{thm:six}
(every listed code has $h(\col)\le2$). Class C1 with
$\operatorname{ord}(\col)$ odd: depth $6$ decided instance by instance
($\star$: depth-$6$ schedule found by the solver; $\dagger$: depth $6$
refuted with a verified DRAT certificate, for the
$\mathbb{Z}_{15}\times\mathbb{Z}_3$ and the $\mathbb{Z}_{45}$ instance
alike). Distances
without ``$\le$'' are exact (Methods); for odd $N$ no element has even
order and every distinct-sums code is also excluded by
Corollary~\ref{cor:twoblock}.}
\label{tab:depth6codes}
\begin{ruledtabular}
\begin{tabular}{cp{0.31\textwidth}p{0.30\textwidth}p{0.27\textwidth}}
$n$ & distinct sums (depth $7$) & class C1, $\operatorname{ord}(\col)$ even (depth $6$) & class C1, $\operatorname{ord}(\col)$ odd \\
\colrule
36  & $[[36,4,6]]$ & $[[36,4,6]]$ & $[[36,4,6]]$ \\
48  & $[[48,8,4]]$, $[[48,4,8]]$ & $[[48,8,4]]$, $[[48,4,8]]$ & $[[48,8,4]]$, $[[48,4,6]]$ \\
54  & $[[54,8,6]]$, $[[54,4,8]]$ & none ($N$ odd) & $[[54,4,8]]$ \\
72  & $[[72,16,4]]$, $[[72,12,6]]$, $[[72,8,8]]$, $[[72,4,\le10]]$ & $[[72,12,4]]$, $[[72,8,6]]$, $[[72,4,8]]$ & $[[72,12,4]]^\star$, $[[72,8,6]]^\star$, $[[72,4,8]]$ \\
90  & $[[90,24,4]]$, $[[90,12,6]]$, $[[90,8,\le10]]$, $[[90,4,\le12]]$ & none ($N$ odd) & $[[90,12,6]]^\dagger$, $[[90,4,\le10]]$ \\
108 & $[[108,24,4]]$, $[[108,16,6]]$, $[[108,8,\le10]]$, $[[108,4,\le12]]$ & $[[108,12,6]]$, $[[108,4,\le12]]$ & $[[108,12,6]]$, $[[108,8,\le8]]$, $[[108,4,\le12]]$ \\
144 & $[[144,32,4]]$, $[[144,24,6]]$, $[[144,16,\le8]]$, $[[144,12,\le12]]$, $[[144,8,\le10]]$, $[[144,4,\le16]]$ & $[[144,24,4]]$, $[[144,16,6]]$, $[[144,12,\le8]]$, $[[144,8,\le10]]$, $[[144,4,\le14]]$ & $[[144,24,4]]$, $[[144,16,6]]$, $[[144,8,\le8]]$, $[[144,4,\le14]]$ \\
\end{tabular}
\end{ruledtabular}
\end{table*}

Up to $n=48$ the depth-six codes match the best depth-seven two-block
codes; from $n=72$ on they trail them by one step in $k$ or in $d$, while
several points coincide. The $[[72,8,6]]$ code over
$\mathbb{Z}_6\times\mathbb{Z}_6$ with $A=1+x^4y^3+x^5y^3$ and
$B=1+xy^5+x^4y^3$ is the smallest example of interest: its only
coincidence is $a_1b_3=a_2b_1=x^4y^3$, so $\col=x^4y^3$ has order $6$ and
$2$-height $0$, and with $\varepsilon(x^iy^j)=j\bmod2$ the class schedule of
Table~\ref{tab:template72} pulls back to a proper depth-$6$ circuit,
verified by $\F$ simulation; the distance $6$ is exact. Compared with the
gross code $[[72,12,6]]$, one layer is traded for four logical qubits at
the same length and distance. The odd-order regime contains such codes
too: the solver found proper depth-$6$ schedules for a $[[72,8,6]]$ code
over $\mathbb{Z}_{36}$ with $\operatorname{ord}(\col)=9$ and for a
$[[72,12,4]]$ code over $\mathbb{Z}_6\times\mathbb{Z}_6$, and refuted depth
$6$ for a $[[90,12,6]]$ code over $\mathbb{Z}_{15}\times\mathbb{Z}_3$.
The distances are code distances; the circuit-level distances of the six-
and seven-layer circuits are not compared here.

\begin{table}[t]
\caption{\textbf{Class schedule of the $[[72,8,6]]$ code.}
$G=\mathbb{Z}_6\times\mathbb{Z}_6$, $A=\{(0,0),(4,3),(5,3)\}$,
$B=\{(0,0),(1,5),(4,3)\}$, collision $a_1+b_3=a_2+b_1$,
$\varepsilon(i,j)=j\bmod2$, key $(0,1,1;0,1,1)$. Pulled back along
$\varepsilon$ it is a proper depth-$6$ schedule, verified by $\F$
simulation.}
\label{tab:template72}
\begin{ruledtabular}
\begin{tabular}{llccc}
class & edge type & label 1 & label 2 & label 3 \\
\colrule
$c=0$ & $X$-check, $L$-edges & 3 & 6 & 1 \\
      & $X$-check, $R$-edges & 2 & 5 & 4 \\
      & $Z$-check, $L$-edges & 5 & 4 & 2 \\
      & $Z$-check, $R$-edges & 3 & 1 & 6 \\
\colrule
$c=1$ & $X$-check, $L$-edges & 5 & 4 & 2 \\
      & $X$-check, $R$-edges & 3 & 6 & 1 \\
      & $Z$-check, $L$-edges & 3 & 1 & 6 \\
      & $Z$-check, $R$-edges & 2 & 4 & 5 \\
\end{tabular}
\end{ruledtabular}
\end{table}

\section{Discussion}
\label{sec:discussion}

Theorems~\ref{thm:main} and~\ref{thm:six} give two exact optima in one
model, seven and six, separated by a single repeated sum. The seven-layer
cycle is optimal for every BB code with distinct products and for every
two-block code over an abelian group with weight-three polynomials and
two-qubit overlaps; the proof is finite combinatorics on the Tanner graph
and applies to codes not yet constructed. A single collision, which changes
neither check weight nor qubit degree, creates four-qubit overlaps in which
the parity constraint is satisfied by cancellation, and when the collision
element has even order and $2$-height at most $8$ the template attains the
degree bound. The algebra of $A$ and $B$, not their weights, decides the
single-round depth. The freedom is compatible with distance: the depth-six
codes of Table~\ref{tab:depth6codes} match the best depth-seven two-block
codes up to $n=48$ and trail them by one step in $k$ or $d$ at
$n=72$--$144$, with $[[72,8,6]]$, $[[108,12,6]]$ and $[[144,16,6]]$ as
representatives. Whether the gap persists in an exhaustive search at
$n=144$, and how the six- and seven-layer circuits compare in circuit-level
distance and logical error rate, remain to be determined; the optimal
depth-seven circuit need not be the best fault-tolerant one, and the two
questions are separate.

The general obstructions reach beyond the family. Theorem~\ref{thm:mod4}
applies to any CSS code whose data qubits saturate the degree bound and
whose $X$-checks have weight $\Delta$, with hypergraph products of
$(3,3)$-regular codes of odd length as examples
(Corollary~\ref{cor:regular}); Theorem~\ref{thm:twisted} applies to any
two-block code, and Theorem~\ref{thm:nomixed} explains why symmetric
constructions stop at $\Delta+1$ for two odd weights. All three hold with
every Tanner-graph edge available as a coupler, unlike the
connectivity-driven bounds of Delfosse \emph{et al.}~\cite{Delfosse2021}.
Their scope is the single-round model: amortised depth $\Delta$ by
overlapping rounds~\cite{Strikis2026,Yoder2025}, morphing
circuits~\cite{ShawTerhal2025,ShawTerhal2026} and flag
schemes~\cite{ChamberlandBeverland2018,ChaoReichardt2018} change the
model, and lower bounds there are open. Depth does not determine fault
tolerance; schedules of equal depth differ in circuit distance, computed
with tools such as \textsc{Stim}~\cite{Gidney2021}, and the residual-error
metrics of Ref.~\cite{Strikis2026} are the instrument for choosing among
them.

Four problems remain open. The rigidity step uses that a depth-$6$
schedule has four interior layers; for $\wA=\wB=5$ Lemma~\ref{lem:parity}
still applies, since $\Delta=10\equiv2\pmod 4$, but the cases
$(a,b)=(1,10),(2,9)$ resist the layer-by-layer argument of
Lemma~\ref{lem:rigidity}, and we state the general odd case with
$\wA+\wB\equiv2\pmod 4$ as a conjecture supported by
Theorems~\ref{thm:mod4}, \ref{thm:nomixed}, \ref{thm:S1}
and~\ref{thm:S2}. When $\Delta\equiv0\pmod 4$, as for $(3,5)$, the parity
mechanism gives no contradiction and we have no evidence either way.
Between the necessary condition of Theorem~\ref{thm:twisted} and the
sufficient condition of Theorem~\ref{thm:six} lie the even-order collision
elements of $2$-height above $8$, where the conjecture is the statement
about the family~\eqref{eq:family}, and the odd-order elements, where the
regime is mapped but not characterised. Every member of the family admits
a class schedule with the half-period symmetry
$T(c+q/2,e)=\sigma_{\mathrm{blk}(e)}(T(c,e))$ for two fixed involutions of
the six layers, for all $q\le512$, but the symmetry does not extend to
$c+q/4$ and the layers on the classes $c<q/2$ depend on $c$ modulo
$2^{j-1}$ in an essential way for most $b$ (Supplementary Section~S5); an
analytic template has to prescribe those. For CSS codes without the
two-block structure,
attainability of the degree bound remains an instance-by-instance question
that the pipeline of Section~\ref{sec:methods} can answer but the theory
does not yet predict.

\section{Methods}
\label{sec:methods}

\paragraph*{Proof of Theorem~\ref{thm:mod4}.}
Fix a component with $n$ data qubits and $m_X$ $X$-checks; since every
$X$-check has weight $\Delta$, $\sum_vr_v=m_X\Delta$. At a data qubit $v$
with $X$-layers $P_v$, the layers above $i\in P_v$ are $Z$-layers except
for the $X$-layers above $i$, so the number of $X$-first corners at $v$ is
$N_v=\sum_{i\in P_v}(\Delta-i)-\binom{r_v}{2}$. Summing over $v$, the first
term is $\sum_e(\Delta-\pi(e))$ over all $X$-edges, which equals
$m_X\Delta(\Delta-1)/2$ because every $X$-check has one edge in each
layer. Hence
\begin{align}
\sum_vN_v
&=\frac{\Delta-1}{2}\sum_vr_v-\sum_v\frac{r_v(r_v-1)}{2}\nonumber\\
&=\frac12\sum_vr_v(\Delta-r_v)=\frac12\sum_vr_vs_v ,
\end{align}
for every valid depth-$\Delta$ layering. Each corner belongs to one
overlapping pair and properness makes each pair's count even, so
$\sum_vr_vs_v\equiv0\pmod 4$. \hfill$\square$

\paragraph*{Proof of Corollary~\ref{cor:twoblock}.}
$X_g$ shares a data qubit with $Z_{g'}$ iff $g'-g\in\Sigma:=\{a_s+b_t\}$,
so the $X$-checks reachable from $X_g$ are $X_{g+H'}$ with
$H'=\langle\sigma-\sigma':\sigma,\sigma'\in\Sigma\rangle$; $H'=H$ because
$a_s-a_{s'}=(a_s+b_t)-(a_{s'}+b_t)$ and every $\sigma-\sigma'$ is the sum
of an $a$-difference and a $b$-difference. The $L$-qubits of $X_{g+H}$ are
$L_{g+a_s+H}$, a single coset since $a_s-a_{s'}\in H$; likewise for the
$R$-qubits $R_{g+b_t+H}$ and the $Z$-checks $Z_{g+a_s+b_t+H}$. Every data
qubit has $r_vs_v=\wA\wB$, so the sum over the component is
$2|H|\wA\wB\equiv2\pmod4$ iff $|H|\wA\wB$ is odd, and
Theorem~\ref{thm:mod4} applies. \hfill$\square$

\paragraph*{Proof of Lemma~\ref{lem:support}.}
Let $x_i$ be the number of $X$-edges in layer $i$ on the $L$ block,
$\Lambda=\{i:x_i>0\}$, $a=\min\Lambda$, $b=\max\Lambda$. In layer $i$ the
$N$ $L$-qubits carry $x_i$ $X$-edges and $N-x_i$ $Z$-edges; the $N-x_i$
$X$-checks with no $L$-edge in layer $i$ have their layer-$i$ edge on the
$R$ block, so the $R$ block carries $N-x_i$ $X$-edges and $x_i$ $Z$-edges
there. Hence $Z$-edges on the $R$ block also lie in $[a,b]$, and every
$X$-check has $V_g\supseteq[1,a-1]\cup[b+1,\Delta]$. We use two properties
of $\Phi$ (Supplementary Section~S2): (P1)~the $\wB$ corners at an
$L$-edge of $X_g$ are mapped, one each, to the $\wB$ $R$-edges of $X_g$,
with equal orientations; (P2)~symmetrically, the $\wB$ corners at an
$R$-edge of a $Z$-check are mapped from its $\wB$ $L$-edges, one from
each.

Take $X_g$ with an $L$-edge in layer $b$. All layers above $b$ at its
$L$-qubit are $Z$-layers, so this edge has exactly $\Delta-b$ $X$-first
corners. By (P1), an $R$-edge of $X_g$ in a layer $v<a$ has its partner
$Z$-edge in $[a,b]$, orientation $1$, and there are $a-1$ such edges; if
$a\in V_g$, the $R$-qubit of the layer-$a$ edge has its $X$-edge in layer
$a$, so the partner lies above $a$, orientation $1$. Hence
$\Delta-b\ge a-1+\mathbb{1}[a\in V_g]$. Take $Z_{g'}$ with an $R$-edge in
layer $a$: it has $a-1$ $X$-first corners, and (P2) gives
$a-1\ge\Delta-b+\mathbb{1}[b\in U_{g'}]$ with $U_{g'}$ the $L$-layer set of
$Z_{g'}$. Adding forces both indicators to vanish and both inequalities to
be equalities: $a+b=\Delta+1$, and $a\in P_g$. Layer reversal
$i\mapsto\Delta+1-i$ preserves properness and fixes $(a,b)$, giving the
converse. \hfill$\square$

\paragraph*{Reduction of the template keys.}
Fix $q=2^j$ and a key with $\colbar=q/2$. Relabel so that the collision is
$\{(1,1),(2,2)\}$ and normalise $\alpha_1=\beta_1=0$ (Supplementary
Section~S5); the collision equation then forces $\alpha_2=\beta_2=q/2$, so
the key is $\kappa_q(a,b)=(0,q/2,a;0,q/2,b)$ with $a=\varepsilon(a_3)$,
$b=\varepsilon(b_3)$: $q^2$ keys per modulus. Three symmetries of the class
system identify further keys: $(a,b)\sim(b,a)$ ($A\leftrightarrow B$ with
$L\leftrightarrow R$), $(a,b)\sim(\omega a,\omega b)$ for every unit
$\omega$ of $\Zq$ (with $T'(c)=T(\omega^{-1}c)$), and
$(a,b)\sim(a+q/2,b+q/2)$ (swap labels $1\leftrightarrow2$ and
renormalise). Moreover the constraints (i)--(iv) only connect classes that
differ by key elements, so the class system of $\kappa_q(a,b)$ is a
disjoint union of copies of the class system of modulus $q/g$ with key
$\kappa_{q/g}(a/g,b/g)$, $g=\gcd(q/2,a,b)$, again with $\colbar$ equal to
half the modulus; it therefore suffices to treat keys with $\gcd(a,b)$
odd, which after a swap and a scaling have $a=1$, the
family~\eqref{eq:family}. The representatives under the three symmetries
number $2,5,11,21,39,73,139,269,527$ for $q=2,4,\dots,512$. For $q\le32$
every representative was decided satisfiable by SAT without further
assumptions ($39$ representatives at $q=32$ in $17$ minutes); for
$q\le512$ every
member of the family~\eqref{eq:family} was decided satisfiable under the
ansatz $T(c+q/2,e)=\sigma_{\mathrm{blk}(e)}(T(c,e))$ with
$\sigma_L=(13)(24)(56)$ on $L$-edges and $\sigma_R=(12)(36)(45)$ on
$R$-edges, whose solutions are class schedules and which makes the
instances easy (all $512$ keys at $q=512$ in twelve minutes); Supplementary
Section~S5 records how the symmetry was found and which stronger forms
fail.

\paragraph*{Class schedules with degenerate keys.}
For a key with repeated entries, $\alpha_s=\alpha_{s'}$ or
$\beta_t=\beta_{t'}$ (unavoidable at $q=2$ and the case for $648$ of the
$864$ admissible keys at $q=4$), the quotient system $\mathrm{QC}(\varepsilon(A),
\varepsilon(B))$ over $\Zq$ is not a two-block code with three-element
sets, so each stored class schedule is checked in two ways: directly
against (i)--(iv), and by pulling it back along $g\mapsto g\bmod q$ to a
\emph{realising} class-C1 code, with lifts $a_s=\alpha_s+qx_s$,
$b_t=\beta_t+qy_t$ in $\mathbb{Z}_{qM}$ chosen so that the designated
collision is the only coincidence, followed by the $\F$ circuit
simulation. All $72$ keys at $q=2$ ($M=7$, codes over $\mathbb{Z}_{14}$)
and all $864$ keys at $q=4$ (codes over $\mathbb{Z}_{28}$) pass.

\paragraph*{Table~\ref{tab:depth7main} and the published cycle.}
In our notation, rounds $1$--$7$ of the cycle of Ref.~\cite{Bravyi2024}
(Table~5 there; $A_s\leftrightarrow$ label $s$, $B_t\leftrightarrow$ label
$t$) are $X$ $L$-edges $(6,2,7)$, $X$ $R$-edges $(4,3,5)$, $Z$ $L$-edges
$(3,4,5)$ and $Z$ $R$-edges $(1,6,2)$. This is a proper
translation-invariant schedule of depth $7$ with
$o_L=o_R=\mathbb{1}[s=2]$, verified by $\F$ simulation on $[[72,12,6]]$ and
$[[144,12,12]]$, whereas Table~\ref{tab:depth7main} has
$o_L=o_R=\mathbb{1}[s\le2]$. The two are not equivalent under label
permutations, layer reversal, $X$/$Z$ duality and block swap, a group of
order $288$ acting on labelled schedules: the $8496$ translation-invariant
proper depth-$7$ schedules fall into $41$ orbits under it, and the two lie
in different orbits. Both attain the minimum of Corollary~\ref{cor:bb}.

\paragraph*{Code search and distance estimation.}
For each group $G=\mathbb{Z}_\ell\times\mathbb{Z}_m$ of order $N$ the pairs
$A=\{0,a_2,a_3\}$, $B=\{0,b_2,b_3\}$ were enumerated (exhaustively for
$N\le36$; $4\times10^5$ random pairs per group for $N\in\{45,54,72\}$),
classified by the multiset of sums, and $k=n-\operatorname{rank}H_X-
\operatorname{rank}H_Z$ was computed over $\F$. Codes with $k>0$ were
reduced to orbit representatives under $\mathrm{Aut}(G)$, translations of
$A$ and $B$, and $(A,B)\mapsto(B,A)$. For each representative an upper
bound on $d$ ($=d_X=d_Z$ by the duality of Supplementary Section~S2) was
obtained by information-set decoding on a generator of $\ker H_Z$, a
vector being a nontrivial logical operator iff it has odd overlap with
some $Z$-logical representative (Prange with Lee--Brickell pairs; $300$
iterations per representative, $2000$ on the Pareto front, $10^4$ on the
codes with explicit schedules). Whenever
$\binom{n}{\lceil d/2\rceil}\le2.5\times10^6$ the bound was confirmed exact
by a meet-in-the-middle enumeration of all vectors of weight at most $d$
in $\ker H_Z$ outside the row space of $H_X$; entries of
Table~\ref{tab:depth6codes} without ``$\le$'' are exact. Depth-$6$
schedules for the listed class-C1 codes of even order were built from class
schedules of modulus $2^{h(\col)+1}$ and verified by $\F$ simulation
(Supplementary Section~S8); the odd-order instances were decided directly
by SAT at depth $6$, in $7$--$600$ seconds each.

\paragraph*{Encoding, solvers and verification.}
A schedule of depth $D$ is encoded with one Boolean variable per
Tanner-edge/layer pair, exactly-one clauses per edge, at-most-one clauses
for every Tanner vertex (data qubits and check ancillas) and layer, and
Eq.~\eqref{eq:parity} as XORs of layer comparisons with Tseitin
auxiliaries, in the manner of SAT-based circuit
synthesis~\cite{Peham2023,TanCong2020}; the encoding accepts any CSS code.
Unsatisfiability is decided by \textsc{kissat}~4.0.4~\cite{Biere2024Kissat}
with DRAT logging and every proof is checked by
\textsc{drat-trim}~\cite{Wetzler2014DRAT,HeuleHuntWetzler2013}, the
standard for computer-assisted refutations~\cite{Heule2016}; satisfiable
instances are decided by \textsc{CaDiCaL}~\cite{Biere2024CaDiCaL} via
PySAT~\cite{Ignatiev2018PySAT}; sweeps use
\textsc{Glucose}~\cite{AudemardSimon2009} under a $60$-second cap, and
capped instances are reported as undetermined, never as unsatisfiable.
Every schedule is verified by propagating the $\F$ register through the
layered CNOT list and comparing with Eq.~\eqref{eq:se-map} at $M_\pi=0$;
the same routine evaluates Eq.~\eqref{eq:parity} on arbitrary layerings.
Class schedules are enumerated by the same encoder on the class system,
stored per admissible key, and their pullbacks verified by the same
simulation. Layer reversal maps proper schedules to proper schedules; an
early encoding broke this symmetry by a clause (Supplementary Section~S7),
and the certificates of the seven codes of Ref.~\cite{Bravyi2024} were
regenerated without it. All runs were single-threaded on one workstation
(Intel Core Ultra 9 275HX, $15.6$~GiB); wall and CPU times agree within a
few per cent, and no cluster resources were used.

\paragraph*{Certificates.}
Eighteen instances carry DRAT certificates: the seven codes of
Ref.~\cite{Bravyi2024} (Table~\ref{tab:certificates}), the smallest and
largest certified distinct-sums instances ($|G|=10$ and $24$),
$G=\mathbb{Z}_6$, $A=B=\{0,2,4\}$, which Theorem~\ref{thm:mod4} excludes
through its six-qubit components and which is the hardest refutation here,
$558$ seconds and a $2.7$~GB proof with the reversal clause retained, and
six class-C1 instances of odd order not covered by a theorem (proofs
$123$--$911$~MB, $16$--$163$~s), and the two odd-order $[[90,12,6]]$ codes
of Supplementary Section~S8 (proofs $2.4$ and $2.8$~GB). As a cross-check, the $17$ certificate
CNFs of the first ten instances were re-decided by \textsc{CaDiCaL} and
\textsc{Glucose} with a $300$-second cap: $16$ are unsatisfiable under
both, and the $\mathbb{Z}_6$ file is decided by \textsc{CaDiCaL} in $233$
seconds but not by \textsc{Glucose} within the cap. All seven codes of
Ref.~\cite{Bravyi2024} are unsatisfiable at depth $6$ without
symmetry-breaking clauses, with solve times between $1.4$ and $6.3$
seconds and verified proofs between $18$ and $68$~MB, and satisfiable at
depth $7$ by Proposition~\ref{prop:seven}. Of the $664$ instances of
Fig.~\ref{fig:sharpness}, $250$ are excluded by Theorems~\ref{thm:mod4}
or~\ref{thm:main}, $80$ of them beyond the solver's cap; $56$ are solver
refutations without proof logging and $6$ carry certificates; $30$ are
open; the remaining $322$ admit schedules verified by simulation, eight of
them found by twisted witnesses after the solver reached its cap.

\begin{table}[t]
\caption{\textbf{Depth-6 refutations of the seven bivariate bicycle codes
of Ref.~\cite{Bravyi2024}.} Group $\mathbb{Z}_\ell\times\mathbb{Z}_m$,
\textsc{kissat} solve time and DRAT proof size, all without
symmetry-breaking clauses and all verified by \textsc{drat-trim}.
Polynomials in Supplementary Table~\ref{tab:instances}.}
\label{tab:certificates}
\begin{ruledtabular}
\begin{tabular}{lccc}
Code & $(\ell,m)$ & Time (s) & Proof (MB) \\
\colrule
$[[72,12,6]]$      & $(6,6)$   & 4.0 & 24 \\
$[[90,8,10]]$      & $(15,3)$  & 4.2 & 41 \\
$[[108,8,10]]$     & $(9,6)$   & 4.8 & 52 \\
$[[144,12,12]]$    & $(12,6)$  & 3.6 & 35 \\
$[[288,12,18]]$    & $(12,12)$ & 1.4 & 18 \\
$[[360,12,\le24]]$ & $(30,6)$  & 6.3 & 68 \\
$[[756,16,\le34]]$ & $(21,18)$ & 3.3 & 46 \\
\end{tabular}
\end{ruledtabular}
\end{table}

\paragraph*{Validation of the model and the lemmas.}
The parity criterion~\eqref{eq:parity} and direct $\F$ simulation agree on
$300$ random valid layerings of each of two instances, the $[[72,12,6]]$
code and the $(\wA,\wB)=(2,3)$ code \texttt{oe23}; this checks the
implementation, the equivalence itself being the derivation of $M_\pi$.
The statements that hold for general $(\wA,\wB)$ with distinct
sums---Lemma~\ref{lem:support}, the corner bijection, $F_L=F_R$, and the
shell statements used in the rigidity proof---hold on all $417$ proper
depth-$\Delta$ schedules of four codes with one even weight,
$(\wA,\wB)\in\{(1,2),(2,1),(2,3),(4,3)\}$, $\Delta\in\{3,3,5,7\}$, of which
$5$ are translation-invariant, $72$ twisted and $340$ without a proper
period; the pair structure of Lemma~\ref{lem:rigidity} cannot be tested
this way, since no proper depth-$6$ schedule of a weight-$(3,3)$ code
exists. The coset identity behind Theorem~\ref{thm:twisted} holds on all
$12{,}931$ cosets formed from $289$ twisted witnesses and all $11{,}412$
admissible template keys, and on the $323$ twisted schedules found on the
satisfiable class-C2 instances, whose $|M|$ takes eleven values between
$2$ and $24$, all even; both subgroup instances are unsatisfiable at depth
$6$, so beyond classes C1 and C2 the theorem rests on its proof alone
(Supplementary Section~S7). Every number in
Sections~\ref{sec:six}--\ref{sec:methods} is anchored to a data file in
the repository.

\section*{Data availability}
All instance definitions, solver logs, DRAT certificates, schedules,
templates and the CSV files underlying every table and figure are
available at \url{https://doi.org/10.5281/zenodo.XXXXXXXX} (release
accompanying this version).

\section*{Code availability}
The SAT encoder, the simulation checker, the enumeration scripts, the
depth-certificate compiler of Supplementary Section~S1, the key-reduction,
code-search and distance-estimation scripts of Sections~\ref{sec:six}
and~\ref{sec:codes}, and a manifest mapping each figure to its data are
available at \url{https://github.com/Guangyi-Leo/A-parity-obstruction-to-degree-optimal-syndrome-extraction-in-bivariate-bicycle-codes}
(archived at
the Zenodo record above); \texttt{run.sh} regenerates every table and
figure.
% The revision scripts and their outputs are in the repository (twoblock/,
% revision/scripts/); run.sh --revision re-runs them.  The Zenodo DOI above
% is still a placeholder: the record has not been deposited yet, so the
% number cannot be filled in here.

\section*{Acknowledgements}
We thank Frederic T.\ Chong, Yongshan Ding, Steve Girvin, Allen Mi, Kate
Smith and Youtao Zhang for helpful discussions. G.L., Y.G., Z.W.\ and
J.L.\ are supported in part by the University of Pittsburgh School of
Computing and Information, the Department of Computer Science, Pitt Cyber,
the Pitt Momentum Fund, the PQI Community Collaboration Awards, the John
C.\ Mascaro Faculty Scholar in Sustainability, NSF award 2535915, Thinking
Machines Lab and Cisco Research. X.Z.\ acknowledges support from the AWS
Center for Quantum Computing and the Fu Foundation School of Engineering
and Applied Science of Columbia University. Z.Z.\ acknowledges NSF CAREER
Award No.~2317471. This research used resources of the Oak Ridge
Leadership Computing Facility, a DOE Office of Science User Facility
supported under Contract DE-AC05-00OR22725.

\section*{Author contributions}
G.L.\ conceived the project, proved the obstruction theorems, the main
theorem and the depth-six theorem, designed the SAT encoding and the
certificate pipeline, ran the computational study, and wrote the
manuscript. Y.G.\ and Z.W.\ contributed to the theoretical analysis and to
the validation of the proofs. X.Z.\ and Z.Z.\ contributed to the analysis
of the circuit model and its hardware assumptions. J.L.\ supervised the
project and revised the manuscript. All authors discussed the results and
reviewed the manuscript.

\section*{Competing interests}
J.L.\ has served as an Associate Editor of npj Quantum Information since
2024. The journal's editorial policies ensure that he was not involved in
the editorial handling or peer review of this manuscript. The other
authors declare no competing interests.

\begin{thebibliography}{99}
\bibitem{BreuckmannEberhardt2021} Breuckmann, N. P. \& Eberhardt, J. N. Quantum low-density parity-check codes. \emph{PRX Quantum} \textbf{2}, 040101 (2021).
\bibitem{Panteleev2022} Panteleev, P. \& Kalachev, G. Asymptotically good quantum and locally testable classical LDPC codes. In \emph{Proc. 54th Annual ACM SIGACT Symposium on Theory of Computing} 375--388 (ACM, 2022).
\bibitem{LeverrierZemor2022} Leverrier, A. \& Z\'emor, G. Quantum Tanner codes. In \emph{Proc. 63rd IEEE Annual Symposium on Foundations of Computer Science} 872--883 (IEEE, 2022).
\bibitem{Dinur2023} Dinur, I., Hsieh, M.-H., Lin, T.-C. \& Vidick, T. Good quantum LDPC codes with linear time decoders. In \emph{Proc. 55th Annual ACM Symposium on Theory of Computing} 905--918 (ACM, 2023).
\bibitem{Bravyi2024} Bravyi, S. \emph{et al.} High-threshold and low-overhead fault-tolerant quantum memory. \emph{Nature} \textbf{627}, 778--782 (2024).
\bibitem{Yoder2025} Yoder, T. J. \emph{et al.} Tour de gross: a modular quantum computer based on bivariate bicycle codes. Preprint at \url{https://arxiv.org/abs/2506.03094} (2025).
\bibitem{Xu2024} Xu, Q. \emph{et al.} Constant-overhead fault-tolerant quantum computation with reconfigurable atom arrays. \emph{Nat. Phys.} \textbf{20}, 1084--1090 (2024).
\bibitem{Voss2024} Voss, L., Xian, S. J., Haug, T. \& Bharti, K. Multivariate bicycle codes. Preprint at \url{https://arxiv.org/abs/2406.19151} (2024).
\bibitem{WangMueller2024} Wang, M. \& Mueller, F. Coprime bivariate bicycle codes and their properties. Preprint at \url{https://arxiv.org/abs/2408.10001} (2024).
\bibitem{Berthusen2025} Berthusen, N. \emph{et al.} Toward a 2D local implementation of quantum low-density parity-check codes. \emph{PRX Quantum} \textbf{6}, 010306 (2025).
\bibitem{LiangChen2026} Liang, Z. \& Chen, Y.-A. Topological subsystem bivariate bicycle codes with four-qubit check operators. Preprint at \url{https://arxiv.org/abs/2605.04151} (2026).
\bibitem{Konig1916} K\H{o}nig, D. \"{U}ber Graphen und ihre Anwendung auf Determinantentheorie und Mengenlehre. \emph{Math. Ann.} \textbf{77}, 453--465 (1916).
\bibitem{TDB2022} Tremblay, M. A., Delfosse, N. \& Beverland, M. E. Constant-overhead quantum error correction with thin planar connectivity. \emph{Phys. Rev. Lett.} \textbf{129}, 050504 (2022).
\bibitem{PMH2008} Patel, K. N., Markov, I. L. \& Hayes, J. P. Optimal synthesis of linear reversible circuits. \emph{Quantum Inf. Comput.} \textbf{8}, 282--294 (2008).
\bibitem{Jiang2020} Jiang, J. \emph{et al.} Optimal space-depth trade-off of CNOT circuits in quantum logic synthesis. In \emph{Proc. 31st Annual ACM-SIAM Symposium on Discrete Algorithms} 213--229 (SIAM, 2020).
\bibitem{MaslovZindorf2022} Maslov, D. \& Zindorf, B. Depth optimization of CZ, CNOT, and Clifford circuits. \emph{IEEE Trans. Quantum Eng.} \textbf{3}, 2500408 (2022).
\bibitem{Kutin2007} Kutin, S. A., Moulton, D. P. \& Smithline, L. M. Computation at a distance. \emph{Chicago J. Theor. Comput. Sci.} \textbf{2007}, 1 (2007).
\bibitem{Bataille2022} Bataille, M. Quantum circuits of CNOT gates: optimization and entanglement. \emph{Quantum Inf. Process.} \textbf{21}, 269 (2022).
\bibitem{DKLP2002} Dennis, E., Kitaev, A., Landahl, A. \& Preskill, J. Topological quantum memory. \emph{J. Math. Phys.} \textbf{43}, 4452--4505 (2002).
\bibitem{Fowler2012} Fowler, A. G., Mariantoni, M., Martinis, J. M. \& Cleland, A. N. Surface codes: towards practical large-scale quantum computation. \emph{Phys. Rev. A} \textbf{86}, 032324 (2012).
\bibitem{TomitaSvore2014} Tomita, Y. \& Svore, K. M. Low-distance surface codes under realistic quantum noise. \emph{Phys. Rev. A} \textbf{90}, 062320 (2014).
\bibitem{Conrad2018} Conrad, J., Chamberland, C., Breuckmann, N. P. \& Terhal, B. M. The small stellated dodecahedron code and friends. \emph{Phil. Trans. R. Soc. A} \textbf{376}, 20170323 (2018).
\bibitem{Zhang2026ASC} Zhang, K. \emph{et al.} Optimal compilation of syndrome extraction circuits for general quantum LDPC codes. Preprint at \url{https://arxiv.org/abs/2603.21499} (2026).
\bibitem{Liu2026AlphaSyndrome} Liu, Y. \emph{et al.} AlphaSyndrome: tackling the syndrome measurement circuit scheduling problem for QEC codes. Preprint at \url{https://arxiv.org/abs/2601.12509} (2026).
\bibitem{Viszlai2026PropHunt} Viszlai, J., Maurya, S., Tannu, S., Martonosi, M. \& Chong, F. T. PropHunt: automated optimization of quantum syndrome measurement circuits. Preprint at \url{https://arxiv.org/abs/2601.17580} (2026).
\bibitem{Nguyen2026} Nguyen, M. T. P., Rimbach-Russ, M. \& Bosco, S. Disassembling qLDPC codes for depth-optimal parity-check circuits. Preprint at \url{https://arxiv.org/abs/2608.19917} (2026).
\bibitem{Bravyi2023arXiv} Bravyi, S. \emph{et al.} High-threshold and low-overhead fault-tolerant quantum memory. Preprint at \url{https://arxiv.org/abs/2308.07915} (2023).
\bibitem{ShawTerhal2025} Shaw, M. H. \& Terhal, B. M. Lowering connectivity requirements for bivariate bicycle codes using morphing circuits. \emph{Phys. Rev. Lett.} \textbf{134}, 090602 (2025).
\bibitem{ShawTerhal2026} Shaw, M. H. \& Terhal, B. M. Optimising quantum error correction using morphing circuits. Preprint at \url{https://arxiv.org/abs/2604.09797} (2026).
\bibitem{ChamberlandBeverland2018} Chamberland, C. \& Beverland, M. E. Flag fault-tolerant error correction with arbitrary distance codes. \emph{Quantum} \textbf{2}, 53 (2018).
\bibitem{ChaoReichardt2018} Chao, R. \& Reichardt, B. W. Quantum error correction with only two extra qubits. \emph{Phys. Rev. Lett.} \textbf{121}, 050502 (2018).
\bibitem{Strikis2026} Strikis, A., Browne, D. E. \& Beverland, M. E. High-performance syndrome extraction circuits for quantum codes. Preprint at \url{https://arxiv.org/abs/2603.05481} (2026).
\bibitem{KovalevPryadko2013} Kovalev, A. A. \& Pryadko, L. P. Quantum Kronecker sum-product low-density parity-check codes with finite rate. \emph{Phys. Rev. A} \textbf{88}, 012311 (2013).
\bibitem{PanteleevKalachev2021} Panteleev, P. \& Kalachev, G. Degenerate quantum LDPC codes with good finite length performance. \emph{Quantum} \textbf{5}, 585 (2021).
\bibitem{LinPryadko2024} Lin, H.-K. \& Pryadko, L. P. Quantum two-block group algebra codes. \emph{Phys. Rev. A} \textbf{109}, 022407 (2024).
\bibitem{WangLinPryadko2023} Wang, R., Lin, H.-K. \& Pryadko, L. P. Abelian and non-Abelian quantum two-block codes. In \emph{Proc. 12th International Symposium on Topics in Coding} (IEEE, 2023).
\bibitem{Delfosse2021} Delfosse, N., Beverland, M. E. \& Tremblay, M. A. Bounds on stabilizer measurement circuits and obstructions to local implementations of quantum LDPC codes. Preprint at \url{https://arxiv.org/abs/2109.14599} (2021).
\bibitem{Gidney2021} Gidney, C. Stim: a fast stabilizer circuit simulator. \emph{Quantum} \textbf{5}, 497 (2021).
\bibitem{Peham2023} Peham, T., Brandl, N., Kueng, R., Wille, R. \& Schmid, L. Depth-optimal synthesis of Clifford circuits with SAT solvers. In \emph{Proc. IEEE International Conference on Quantum Computing and Engineering} 802--813 (IEEE, 2023).
\bibitem{TanCong2020} Tan, B. \& Cong, J. Optimal layout synthesis for quantum computing. In \emph{Proc. 39th International Conference on Computer-Aided Design} 1--9 (ACM, 2020).
\bibitem{Biere2024Kissat} Biere, A., Faller, T., Fazekas, K., Fleury, M., Froleyks, N. \& Pollitt, F. CaDiCaL, Gimsatul, IsaSAT and Kissat entering the SAT Competition 2024. In \emph{Proc. SAT Competition 2024 -- Solver, Benchmark and Proof Checker Descriptions} 8--10 (University of Helsinki, 2024).
\bibitem{Wetzler2014DRAT} Wetzler, N., Heule, M. J. H. \& Hunt, W. A. DRAT-trim: efficient checking and trimming using expressive clausal proofs. In \emph{Theory and Applications of Satisfiability Testing -- SAT 2014}, Lecture Notes in Computer Science Vol.~8561, 422--429 (Springer, 2014).
\bibitem{HeuleHuntWetzler2013} Heule, M. J. H., Hunt, W. A. \& Wetzler, N. Trimming while checking clausal proofs. In \emph{Proc. Formal Methods in Computer-Aided Design} 181--188 (IEEE, 2013).
\bibitem{Heule2016} Heule, M. J. H., Kullmann, O. \& Marek, V. W. Solving and verifying the Boolean Pythagorean triples problem via cube-and-conquer. In \emph{Theory and Applications of Satisfiability Testing -- SAT 2016}, Lecture Notes in Computer Science Vol.~9710, 228--245 (Springer, 2016).
\bibitem{Biere2024CaDiCaL} Biere, A. \emph{et al.} CaDiCaL 2.0. In \emph{Computer Aided Verification -- CAV 2024}, Lecture Notes in Computer Science Vol.~14681, 133--152 (Springer, 2024).
\bibitem{Ignatiev2018PySAT} Ignatiev, A., Morgado, A. \& Marques-Silva, J. PySAT: a Python toolkit for prototyping with SAT oracles. In \emph{Theory and Applications of Satisfiability Testing -- SAT 2018}, Lecture Notes in Computer Science Vol.~10929, 428--437 (Springer, 2018).
\bibitem{AudemardSimon2009} Audemard, G. \& Simon, L. Predicting learnt clauses quality in modern SAT solvers. In \emph{Proc. 21st International Joint Conference on Artificial Intelligence} 399--404 (AAAI Press, 2009).
% ---- cited only in the Supplementary Information ----
\bibitem{DehaeneDeMoor2003} Dehaene, J. \& De Moor, B. Clifford group, stabilizer states, and linear and quadratic operations over GF(2). \emph{Phys. Rev. A} \textbf{68}, 042318 (2003).
\bibitem{Gottesman1998} Gottesman, D. The Heisenberg representation of quantum computers. Preprint at \url{https://arxiv.org/abs/quant-ph/9807006} (1998).
\bibitem{Aaronson2004} Aaronson, S. \& Gottesman, D. Improved simulation of stabilizer circuits. \emph{Phys. Rev. A} \textbf{70}, 052328 (2004).
\bibitem{ColeOstSchirra2001} Cole, R., Ost, K. \& Schirra, S. Edge-coloring bipartite multigraphs in $O(E\log D)$ time. \emph{Combinatorica} \textbf{21}, 5--12 (2001).
\bibitem{GabowKariv1982} Gabow, H. N. \& Kariv, O. Algorithms for edge coloring bipartite graphs and multigraphs. \emph{SIAM J. Comput.} \textbf{11}, 117--129 (1982).
\bibitem{MooreNilsson2001} Moore, C. \& Nilsson, M. Parallel quantum computation and quantum codes. \emph{SIAM J. Comput.} \textbf{31}, 799--815 (2001).
\bibitem{deBrugiere2021} de Brugi\`ere, T. G., Baboulin, M., Valiron, B., Martiel, S. \& Allouche, C. Gaussian elimination versus greedy methods for the synthesis of linear reversible circuits. \emph{ACM Trans. Quantum Comput.} \textbf{2}, 11 (2021).
\bibitem{Jaques2023} Jaques, S. \& Rattew, A. G. QRAM: a survey and critique. \emph{Quantum} \textbf{9}, 1922 (2025).
\end{thebibliography}

% =====================================================================
% Supplementary Information (same document, shared bibliography)
% =====================================================================
\clearpage
\onecolumngrid
\begin{center}\Large\bfseries Supplementary Information\end{center}
\twocolumngrid
\raggedbottom
\providecommand{\SIwidefigurewidth}{\textwidth}
\providecommand{\SIcolumnfigurewidth}{\columnwidth}
\setcounter{topnumber}{3}
\setcounter{bottomnumber}{2}
\setcounter{totalnumber}{5}
\setcounter{dbltopnumber}{2}
\renewcommand{\topfraction}{0.92}
\renewcommand{\bottomfraction}{0.85}
\renewcommand{\textfraction}{0.06}
\renewcommand{\floatpagefraction}{0.75}
\renewcommand{\dbltopfraction}{0.92}
\renewcommand{\dblfloatpagefraction}{0.72}
\setlength{\textfloatsep}{10pt plus 3pt minus 3pt}
\setlength{\floatsep}{9pt plus 3pt minus 2pt}
\setlength{\intextsep}{9pt plus 3pt minus 2pt}
\setcounter{section}{0}
\renewcommand{\thesection}{S\arabic{section}}
\setcounter{figure}{0}\renewcommand{\thefigure}{S\arabic{figure}}
\setcounter{table}{0}\renewcommand{\thetable}{S\arabic{table}}
\setcounter{equation}{0}\renewcommand{\theequation}{S\arabic{equation}}
\setcounter{theorem}{0}\renewcommand{\thetheorem}{S\arabic{theorem}}

Numbered statements of the main text are cited by their main-text
numbers; the labels S1--S3 are reserved for the three theorems proved only
here, and the unnumbered lemma and proposition of Section~S1 are
background. Notation is that of Section~\ref{sec:model}; in a
depth-$\Delta$ schedule every data qubit and every check has exactly one
edge in each layer. Fig.~\ref{fig:fullmap} is the full dependency map.

\begin{figure*}[!t]
\centering
\includegraphics[width=\SIwidefigurewidth]{figS1_proof_map.pdf}
\caption{\textbf{Full logical structure of the results.} Solid boxes are
proved statements, dashed boxes computational evidence, dotted boxes open
problems. The three general obstructions follow from $M_\pi=0$ alone. The
main chain runs through Lemma~\ref{lem:support} to
Lemmas~\ref{lem:rigidity} and~\ref{lem:parity}; together with
Proposition~\ref{prop:seven} it gives Corollary~\ref{cor:bb}. The
depth-six chain runs from the template enumeration and the reduction of
the keys through Lemma~\ref{lem:lift} and the $2$-height
(Lemma~\ref{lem:height}) to Theorem~\ref{thm:six}.}
\label{fig:fullmap}
\end{figure*}

\section{The depth certificate for commuting gate sets}
\label{si:compiler}

The degree bound of Section~\ref{sec:model} and the K\H{o}nig colouring
that attains it are instances of a depth-exact compiler for affine maps;
this section fixes what the certificate is exact for.

\begin{lemmaSI}
A Clifford circuit $C$ on $n+k$ qubits that maps every computational-basis
input $|\mathbf{x}\rangle|0\rangle^{\otimes k}$ to a computational-basis
output $|g(\mathbf{x})\rangle|h(\mathbf{x})\rangle$ acts by affine maps
$g,h$ over $\F$~\cite{DehaeneDeMoor2003,Gottesman1998,Aaronson2004}.
Conversely, every affine $f(\mathbf{x})=A\mathbf{x}+\mathbf{b}$ is
implemented on $n+m$ qubits by
$C_f=\prod_j\bigl[\bigl(\prod_{k:A_{j,k}=1}\mathrm{CNOT}(x_k,d_j)\bigr)
X_{d_j}^{b_j}\bigr]$.
\end{lemmaSI}

\begin{proof}
Write $CZ_iC^\dagger=(-1)^{\sigma_i}P_i$. A Pauli word with an $X$ or $Y$
factor cannot stabilise a basis state, so $P_i=Z^{\mathbf{v}_i}$; if
$\sum_iu_i\mathbf{v}_i=0$ with $\mathbf{u}\neq0$ then
$\prod_{u_i=1}Z_i=\pm I$, which is false, so
$V=(\mathbf{v}_1\cdots\mathbf{v}_{n+k})$ is invertible. The output
$\mathbf{z}$ satisfies $V^{\mathsf T}\mathbf{z}=(\mathbf{x},0)+\bm{\sigma}$,
hence is affine in $\mathbf{x}$. The converse is a direct evaluation.
\end{proof}

The gates of $C_f$ commute pairwise, so compilation is a scheduling
problem. Read $(A,\mathbf{b})$ as a bipartite multigraph $G'_A$ with one
edge per nonzero $A_{j,k}$ and a pendant edge per $b_j=1$; a layering is an
edge colouring of $G'_A$.

\begin{propSI}
Among all schedules of the gate multiset of $C_f$ the minimum depth is
$D^\star=\max\{\max_kc_k,\max_j(w_j+b_j)\}=\Delta(G'_A)$, attained by any
edge colouring of $G'_A$ with $D^\star$ colours~\cite{Konig1916}, computable
in $O(|E'|\log D^\star)$ time~\cite{ColeOstSchirra2001,GabowKariv1982}.
\end{propSI}

$D^\star$ is exact for this gate set and not a lower bound over all CNOT
circuits: with CNOTs inside the data register an all-ones column is
realised in depth $\lceil\log_2(m+1)\rceil$~\cite{MooreNilsson2001,
Jiang2020,MaslovZindorf2022,deBrugiere2021}. It depends only on the
degree sequence of $G'_A$, not on the rank of $A$ or on hardware
connectivity (for the relation of such compiled lookups to quantum memory
see Ref.~\cite{Jaques2023}). The two halves of an extraction round, $A=H_X^{\mathsf T}$ and $A=H_Z$, are
each scheduled exactly by the proposition; interleaving them is an edge
colouring of the Tanner graph, attained with $\Delta$ colours, but the
halves do not commute and the interleaved circuit implements
Eq.~\eqref{eq:se-map} with the cross term $M_\pi$.

\section{Two-block codes and the corner bijection}
\label{si:corners}

\paragraph*{Proof of Eq.~\eqref{eq:heisenberg}.}
Under $\mathrm{CNOT}(c,t)$,
\[
X_c\mapsto X_cX_t,
\qquad
Z_t\mapsto Z_cZ_t,
\]
while $X_t$ and $Z_c$ are fixed. We propagate the Pauli operators
backwards through the scheduled circuit.

First consider $Z_{z_j}$. At a $Z$-check CNOT
$\mathrm{CNOT}(q,z_j)$, it acquires a factor $Z_q$. This factor is
subsequently carried onto $Z_{a_i}$ when the backward propagation reaches
an earlier $X$-check CNOT $\mathrm{CNOT}(a_i,q)$. Thus a shared data qubit
$q\in\operatorname{supp}X_i\cap\operatorname{supp}Z_j$ contributes
$Z_{a_i}$ precisely when its $X$-edge precedes its $Z$-edge. By
definition, the number of such shared qubits is $(M_\pi)_{ji}$ modulo
$2$. Therefore
\[
U^\dagger Z_{z_j}U
=
Z_{z_j}\,\overline{Z}_j
\prod_i Z_{a_i}^{(M_\pi)_{ji}}.
\]

Now consider $X_{a_i}$. At an $X$-check CNOT
$\mathrm{CNOT}(a_i,q)$, it acquires a factor $X_q$. Under backward
propagation, this factor is carried onto $X_{z_j}$ when the circuit
encounters an earlier $Z$-check CNOT $\mathrm{CNOT}(q,z_j)$. Hence a
shared data qubit contributes $X_{z_j}$ precisely when its $Z$-edge
precedes its $X$-edge.

Let
\[
r_{ij}
=
\bigl|
\operatorname{supp}X_i
\cap
\operatorname{supp}Z_j
\bigr|.
\]
Because the two incident CNOTs on a shared data qubit cannot occupy the
same layer, the number of $Z$-first shared qubits is
\[
r_{ij}-(M_\pi)_{ji}.
\]
The CSS commutation condition
\[
H_XH_Z^{\mathsf T}=0
\qquad\text{over }\F
\]
implies that every $X_i$--$Z_j$ overlap has even cardinality:
\[
r_{ij}\equiv0\pmod2.
\]
Consequently,
\[
r_{ij}-(M_\pi)_{ji}
\equiv
(M_\pi)_{ji}
\pmod2.
\]
It follows that
\[
U^\dagger X_{a_i}U
=
X_{a_i}\,\overline{X}_i
\prod_j X_{z_j}^{(M_\pi)_{ji}}.
\]
Thus Eq.~\eqref{eq:heisenberg} has the stated form. The two evolved
operators commute for every $i,j$, as required. \hfill$\square$

\paragraph*{Conventions.}
For BB codes, $G=\mathbb{Z}_\ell\times\mathbb{Z}_m$ and $A,B$ are the
exponent sets of the polynomials of Ref.~\cite{Bravyi2024}, a monomial
$x^iy^j$ being the element $(i,j)$. Table~\ref{tab:instances} lists the
seven codes of Ref.~\cite{Bravyi2024}; all have distinct sums.
Equivalently, the integer matrix $H_XH_Z^{\mathsf T}$ has entries in
$\{0,2\}$, since its $(g,g')$ entry equals $2S(g'-g)$.

\begin{table*}[!t]
\caption{\textbf{The seven bivariate bicycle codes of
Ref.~\cite{Bravyi2024}.} Group, exponent sets $A$ and $B$ with the monomial
$x^iy^j$ written $(i,j)$, whether the nine sums are pairwise distinct, and
the parity of $\ell m$ relevant to Theorem~\ref{thm:mod4}.}
\label{tab:instances}
\begin{ruledtabular}
\begin{tabular}{lcllcc}
Code & $(\ell,m)$ & $A$ & $B$ & distinct sums & $\ell m$ \\
\colrule
$[[72,12,6]]$
  & $(6,6)$
  & $\{(3,0),(0,1),(0,2)\}$
  & $\{(0,3),(1,0),(2,0)\}$
  & yes & even \\
$[[90,8,10]]$
  & $(15,3)$
  & $\{(9,0),(0,1),(0,2)\}$
  & $\{(0,0),(2,0),(7,0)\}$
  & yes & odd \\
$[[108,8,10]]$
  & $(9,6)$
  & $\{(3,0),(0,1),(0,2)\}$
  & $\{(0,3),(1,0),(2,0)\}$
  & yes & even \\
$[[144,12,12]]$
  & $(12,6)$
  & $\{(3,0),(0,1),(0,2)\}$
  & $\{(0,3),(1,0),(2,0)\}$
  & yes & even \\
$[[288,12,18]]$
  & $(12,12)$
  & $\{(3,0),(0,2),(0,7)\}$
  & $\{(0,3),(1,0),(2,0)\}$
  & yes & even \\
$[[360,12,\le24]]$
  & $(30,6)$
  & $\{(9,0),(0,1),(0,2)\}$
  & $\{(0,3),(25,0),(26,0)\}$
  & yes & even \\
$[[756,16,\le34]]$
  & $(21,18)$
  & $\{(3,0),(0,10),(0,17)\}$
  & $\{(0,5),(3,0),(19,0)\}$
  & yes & even \\
\end{tabular}
\end{ruledtabular}
\end{table*}

\paragraph*{Overlaps.}
$X_g$ acts on $L_{g+a_s}$ and $R_{g+b_t}$, while $Z_{g'}$ acts on
$L_{g'-b_t}$ and $R_{g'-a_s}$. They share $L_h$ if and only if
\[
h=g+a_s=g'-b_t,
\]
and share $R_r$ if and only if
\[
r=g+b_t=g'-a_s,
\]
both of which give
\[
g'-g=a_s+b_t.
\]
Each label pair satisfying $a_s+b_t=g'-g$ contributes one shared
$L$-qubit, $L_{g+a_s}$, where the $X$-edge has label $s$ and the
$Z$-edge has label $t$, and one shared $R$-qubit, $R_{g+b_t}$, where the
labels are $t$ and $s$, respectively.

\begin{lemmaSI}[corner bijection]
Assume distinct sums. The map $\Phi$ sending the $L$-corner
\[
(L_{g+a_s},e^{X_g}_s,e^{Z_{g'}}_t)
\]
to the $R$-corner
\[
(R_{g+b_t},e^{X_g}_t,e^{Z_{g'}}_s),
\qquad
g'=g+a_s+b_t,
\]
is a bijection from $L$-corners to $R$-corners, and a schedule is proper
if and only if $\Phi$ preserves orientation.

\emph{(P1)} The $\wB$ corners at the $L$-edge of label $s$ of $X_g$ are
mapped to the $\wB$ $R$-edges of $X_g$, one to each, and the partner
$Z$-edge at the $R$-edge of label $t$ has label $s$.

\emph{(P2)} The $\wB$ corners at the $R$-edge of label $s$ of $Z_{g'}$
(one per $X$-check $X_{g'-a_s-b_t}$ at $R_{g'-a_s}$) are mapped from the
$\wB$ $L$-edges of $Z_{g'}$, one from each.
\end{lemmaSI}

\begin{proof}
An $L$-corner at $L_h$ pairs $X_g$ with
\[
h=g+a_s
\]
and $Z_{g'}$ with
\[
h=g'-b_t.
\]
Hence $g'-g=a_s+b_t$, and distinct sums determine $(s,t)$ from $g'-g$.
The pair shares exactly one corresponding $R$-qubit, namely
$R_{g+b_t}$, with the stated labels. Equation~\eqref{eq:parity} for the
pair says that the two corners have equal orientation. Properties (P1)
and (P2) then follow directly from the labels.
\end{proof}

\paragraph*{Symmetries.}
Layer reversal $i\mapsto\Delta+1-i$ flips every orientation and preserves
properness; it maps $(a,b)$ to $(\Delta+1-b,\Delta+1-a)$, which is $(a,b)$
for proper schedules by Lemma~\ref{lem:support}. The duality
\[
X_g\leftrightarrow Z_{-g},
\qquad
L_h\leftrightarrow R_{-h},
\]
maps $\mathrm{QC}(A,B)$ to itself with $X$ and $Z$ exchanged and
preserves properness.

\section{General obstructions}
\label{si:obstructions}

\paragraph*{Proof of Theorem~\ref{thm:nomixed}.}
Suppose no layer is mixed. Then every $X$-check has the same $L$-layer set
$P$ and $R$-layer set $V=[\Delta]\setminus P$, $|P|=\wA$, $|V|=\wB$; every
$L$-qubit has $X$-layers $P$ and $Z$-layers $V$, every $R$-qubit the
reverse. For any $X$-check $F_L=\#\{(i,v)\in P\times V:i<v\}$ and
$F_R=\#\{(v,i)\in V\times P:v<i\}$, so $F_L+F_R=\wA\wB$, which
Eq.~\eqref{eq:parity} makes even. A translation-invariant schedule has no
mixed layer. \hfill$\square$

\paragraph*{Proof of Theorem~\ref{thm:twisted}.}
Let $\alpha_s=\varepsilon(a_s)$, $\beta_t=\varepsilon(b_t)$. A twisted
schedule assigns to the $X$-check of class $c$ layers $\pi_c(s)$,
$\nu_c(t)$ and to the $Z$-check of class $c$ layers $\mu_c(t)$,
$\kappa_c(s)$. The $L$-qubit of class $\lambda$ has $X$-edges from classes
$\lambda-\alpha_s$ and $Z$-edges from classes $\lambda+\beta_t$; the
$R$-qubit of class $\rho$ has $X$-edges from $\rho-\beta_t$ and $Z$-edges
from $\rho+\alpha_s$. Fix a coset $C$ of $M$ in $\varepsilon(G)$; since the
$\alpha_s$ are congruent modulo $M$, as are the $\beta_t$, the sets
$C+\alpha_s$, $C+\beta_t$, $C+\alpha_s+\beta_t$ are single cosets $C_L$,
$C_R$, $C_Z$ of size $|M|$, and the classes in $C\cup C_L\cup C_R\cup C_Z$
form a closed subsystem. For $c\in C$ let
$F_L(c)=\sum_{s,t}\mathbb{1}[\pi_c(s)<\mu_{c+\alpha_s+\beta_t}(t)]$ and
$F_R(c)=\sum_{s,t}\mathbb{1}[\nu_c(t)<\kappa_{c+\alpha_s+\beta_t}(s)]$;
these are the $F_L,F_R$ of any check of class $c$. Regrouping by qubit
class, at $\lambda\in C_L$ the $\wA$ $X$-edges and $\wB$ $Z$-edges fill
all $\Delta$ layers, so the $X$-first corners there number
$\sum_{i\in P_\lambda}(\Delta-i)-\binom{\wA}{2}$; summing over $C_L$ the
first term is $\sum_{c\in C}\sum_s(\Delta-\pi_c(s))$, and likewise on the
$R$ block with $\binom{\wB}{2}$. Since the $\Delta$ edges of each $X$-check
fill all layers, $\sum_s(\Delta-\pi_c(s))+\sum_t(\Delta-\nu_c(t))=
\binom{\Delta}{2}$, and
\begin{equation}
\begin{aligned}
\sum_{c\in C}\bigl(F_L(c)+F_R(c)\bigr)
&=|M|\Bigl[\tbinom{\Delta}{2}-\tbinom{\wA}{2}-\tbinom{\wB}{2}\Bigr]\\
&=|M|\,\wA\wB .
\end{aligned}
\end{equation}
Each summand is a sum over overlapping pairs of even numbers, hence even.
So $|M|\wA\wB$ is even; for $\wA=\wB=3$, $|M|$ is even, and a subgroup of
$\Zq$ with $q$ odd has odd order. \hfill$\square$

\begin{remarkSI}
With $q=1$ the theorem gives the translation-invariant half of
Theorem~\ref{thm:nomixed}. The proof uses neither distinct sums nor the
class of the code.
\end{remarkSI}

\section{Proof of the main theorem}
\label{si:main}

The proof of Lemma~\ref{lem:support} is in Methods. The lemma says nothing
about layers strictly between $a$ and $b$: a proper depth-$4$ schedule of
$G=\mathbb{Z}_5$, $A=\{0,1\}$, $B=\{0,2\}$ has every $X$-check's $L$-layer
set equal to $\{1,4\}$, with $a+b=5$ and the middle layers unused.
Throughout this section the code has distinct sums.

\begin{theorem}
\label{thm:S1}
Let $\mathrm{QC}(A,B)$ have distinct sums, $|A|=1$ and $|B|=\wB$ odd.
Then no proper schedule of depth $\Delta=\wB+1$ exists; by duality the
same holds for $|B|=1$, $|A|$ odd.
\end{theorem}

\begin{proof}
Let $p_g$, $q_{g'}$, $p_h$, $q_r$ be the layers of the single $L$-edge of
$X_g$, the single $R$-edge of $Z_{g'}$, the single $X$-edge at $L_h$ and
the single $Z$-edge at $R_r$. The $L$-edge of $X_g$ has $\Delta-p_g$
corners, all $X$-first, and by (P1)
$\Delta-p_g=\sum_t\mathbb{1}[v_t<q_{r_t}]$ over its $R$-edges in layers
$[\Delta]\setminus\{p_g\}$; the $R$-edge of $Z_{g'}$ has $q_{g'}-1$
corners, all $X$-first, and by (P2)
$q_{g'}-1=\sum_t\mathbb{1}[p_{h_t}<u_t]$. Layer counting gives
$\#\{g':q_{g'}=i\}=\#\{g:p_g=i\}$, so $p$ and $q$ take values in the same
$[a,b]$ with both ends attained. If $a<b$, take $p_g=b$: $R$-edges in
layers $v<a$ contribute $a-1$, and the $R$-edge in layer $a$ has
$q_r\neq a$, $q_r\ge a$, contributing $1$, so $\Delta-b\ge a$; take
$q_{g'}=a$: $L$-edges in layers $u>b$ contribute $\Delta-b$ and the
$L$-edge in layer $b$ has $p_h<b$, contributing $1$, so
$a-1\ge\Delta-b+1$, a contradiction. If $a=b$, the $Z$-side identity gives
$\Delta=2a-1$, odd; but $\Delta=\wB+1$ is even.
\end{proof}

\begin{theorem}
\label{thm:S2}
Let $\mathrm{QC}(A,B)$ have distinct sums, $|A|=\wA\ge3$ odd and
$|B|=\wB$ odd, and suppose a proper schedule of depth $\Delta=\wA+\wB$ has
$a=(\wB+1)/2$. Then no such schedule exists.
\end{theorem}

\begin{proof}
Here $b=\wA+(\wB+1)/2$, $[a,b]$ has $\wA+1$ layers, and the shells
$S_-=[1,a-1]$, $S_+=[b+1,\Delta]$ have $(\wB-1)/2$ layers each. Each
$X$-check uses $a$ or $b$, hence both by Lemma~\ref{lem:support}:
$P_g=[a,b]\setminus\{c(g)\}$, $V_g=S_-\cup S_+\cup\{c(g)\}$. In layers $a,b$
every $X$-check is on $L$, so every $L$-qubit has $X$-layers
$[a,b]\setminus\{z_h\}$ and every $R$-qubit $Z$-layers
$[a,b]\setminus\{v_r\}$, with the interior $R$-edge of $X_g$ landing where
$v_r=c(g)$; layer counting gives $\#\{h:z_h=i\}=\#\{g:c(g)=i\}$, so the
$z_h$ lie in $[c_{\min},c_{\max}]$. Counting $Z$-layers above each edge,
$F_L=\wA(a-1)+1+W$ with $W=\#\{p\in(a,b)\setminus\{c\}:z_{h_p}>p\}$ and
$F_R=\wA(a-1)+(b-c)$, so $F_L=F_R$ gives $W(g)=b-c(g)-1$. If
$c_{\min}=c_{\max}=c$ then $z_h=c$ for all $h$, $W=c-a-1$ and
$2c=\Delta+1$, impossible for even $\Delta$. Otherwise $c(g)=c_{\max}$
gives $W\ge c_{\min}-a$, hence $c_{\min}+c_{\max}\le\Delta$, and
$c(g)=c_{\min}$ gives $W\le c_{\max}-a-2$, hence
$c_{\min}+c_{\max}\ge\Delta+2$.
\end{proof}

\begin{remarkSI}
The parity of $\wA$ cannot be dropped from Theorem~\ref{thm:S2}: the last
step needs $\Delta$ even. For $G=\mathbb{Z}_{12}$, $A=\{0,1\}$,
$B=\{0,4,8\}$ ($\wA=2$, $\wB=3$, $\Delta=5$) a proper depth-$5$ schedule
with $a=2$, $b=4$ exists, in which every $X$-check has $L$-layers $\{2,4\}$,
so $c(g)=3$ for all $g$ and $2c=6=\Delta+1$ (found by SAT and verified by
$\F$ simulation).
\end{remarkSI}

\begin{theorem}[the case $(2,5)$]
\label{thm:S3}
Let $|A|=|B|=3$ with distinct sums. No proper schedule of depth $6$ has
$(a,b)=(2,5)$.
\end{theorem}

\begin{proof}
In layers $2$ and $5$ every $X$-check has its edge on the $L$ block
($2,5\in P_g$ for all $g$, by Lemma~\ref{lem:support} and $|P_g|=3$ inside
$[2,5]$), so every $L$-qubit has $X$-edges in layers $2$ and $5$ and,
since its $X$-layers lie in $[a,b]=[2,5]$, a third $X$-layer
$c_L\in\{3,4\}$; its $Z$-layers are the complement $\{1,6,7-c_L\}$. On the
$R$ block layers $2$ and $5$ carry no $X$-edge, so every $R$-qubit has
$Z$-edges in layers $2$ and $5$ and a third $Z$-layer $c_R\in\{3,4\}$ (the
$Z$-edges on $R$ lie in $[a,b]$); its $X$-layers are $\{1,6,7-c_R\}$. Fix
$X_g$ with $P_g=\{2,5,c\}$, $V_g=\{1,6,7-c\}$. Its layer-$c$ $L$-edge lies
on an $L$-qubit whose $X$-layers contain $c$, so $c_L=c$ there; its
layer-$(7-c)$ $R$-edge lies on an $R$-qubit whose $X$-layers contain
$7-c$, so $c_R=c$ there. Counting $Z$-layers above each edge,
\begin{align*}
F_L(g)&=|\{6,7-c_L\}\cap(2,6]|+|\{6\}|+|\{6,7-c\}\cap(c,6]|\\
      &=2+1+(1+\mathbb{1}[c=3]),\\
F_R(g)&=|\{2,5,c_R\}|+0+|\{2,5,c\}\cap(7-c,6]|\\
      &=3+0+(1+\mathbb{1}[c=4]),
\end{align*}
using $7-c_L\in\{3,4\}>2$ in the first term, $7-c>c$ iff $c=3$ in the
third, and $5>7-c\ge3>2$, $c>7-c$ iff $c=4$ in the last. So
$F_L(g)=4+\mathbb{1}[c=3]$ and $F_R(g)=4+\mathbb{1}[c=4]$ differ by one
for either value of $c$, contradicting $F_L=F_R$.
\end{proof}

\paragraph*{Proof of Lemma~\ref{lem:rigidity}.}
Let $\wA=\wB=3$, the sums distinct, $\pi$ a proper schedule of depth $6$
with $(a,b)=(1,6)$, and $J=\{2,3,4,5\}$. Let $x_i$ be the number of
$X$-edges in layer $i$ on the $L$ block; as in the proof of
Lemma~\ref{lem:support}, layer $i$ then carries $N-x_i$ $X$-edges and
$x_i$ $Z$-edges on the $R$ block.

\emph{Types.} An $X$-check $g$ with $1\in P_g$ has $6\in P_g$ by
Lemma~\ref{lem:support}; call it an $E$-check and write
$P_g=\{1,6,c_g\}$, $c_g\in J$, $V_g=J\setminus\{c_g\}$. An $X$-check with
$1\notin P_g$ has $6\notin P_g$ by the converse half of the lemma, so
$P_g\subset J$, $|P_g|=3$: call it an $I$-check and write
$P_g=J\setminus\{d_g\}$, $V_g=\{1,6,d_g\}$. Let $E$, $I$ be the two sets
and $e=|E|$. Since $a=1$, some $X$-check has an $L$-edge in layer $1$, so
$e=x_1\ge1$. On the $R$ block layer $1$ carries $x_1=e$ $Z$-edges, one
per $Z$-check that has an $R$-edge in layer $1$; applying
Lemma~\ref{lem:support} to the dual code ($X_g\leftrightarrow Z_{-g}$,
$L\leftrightarrow R$), whose lowest $X$-on-$L$ layer is therefore $1$,
these $Z$-checks also have an $R$-edge in layer $6$. Every $X$-check has
exactly one $L$-edge of each label $s$ and one $R$-edge of each label
$t$; write $s_i(g)$ for the label of the $L$-edge of $X_g$ in layer
$i\in P_g$ and $t_i(g)$ for the label of its $R$-edge in layer
$i\in V_g$.

\emph{(T1)} Let $g\in E$ and let $t$ be any label, $r=R_{g+b_t}$ the
$R$-qubit of the $R$-edge of $X_g$ of label $t$, in layer $\nu_g(t)$.
Then the $Z$-edge at $r$ of label $s_1(g)$, which belongs to
$Z_{g+a_{s_1(g)}+b_t}$, lies in a layer above $\nu_g(t)$, and the
$Z$-edge at $r$ of label $s_6(g)$, belonging to $Z_{g+a_{s_6(g)}+b_t}$,
lies below $\nu_g(t)$. Proof: the $L$-edge of $X_g$ in layer $1$, at
$h=L_{g+a_{s_1(g)}}$, has three corners, with the $Z$-edges at $h$ of
labels $t=1,2,3$; all are $X$-first. By (P1) the corner with label $t$
is mapped by $\Phi$ to the corner at $r=R_{g+b_t}$ formed by the $R$-edge
of $X_g$ of label $t$ and the $Z$-edge at $r$ of label $s_1(g)$, both
edges of the same pair $(X_g,Z_{g+a_{s_1(g)}+b_t})$; properness makes
this corner $X$-first, which is the first claim. The layer-$6$ $L$-edge
has all corners $Z$-first, and the same map gives the second claim.

\emph{(T2)} Let $g\in I$ and $s$ any label, $h=L_{g+a_s}$ the $L$-qubit
of the $L$-edge of $X_g$ of label $s$. Then the $Z$-edge at $h$ of label
$t_1(g)$ lies above that $L$-edge and the $Z$-edge at $h$ of label
$t_6(g)$ lies below it. Proof: the $R$-edge of $X_g$ in layer $1$ has
all three corners $X$-first; by (P1) applied in the reverse direction,
the corner with $Z$-label $s$ there is the image of the corner at $h$
formed by the $L$-edge of label $s$ and the $Z$-edge of label $t_1(g)$,
so that corner is $X$-first. Layer $6$ likewise.

\emph{Step 1: every $L$-qubit with an $X$-edge in layer $5$ has a
$Z$-edge in layer $6$.} Let $h$ carry the layer-$5$ $L$-edge of $X_g$,
of label $s^\ast=s_5(g)$. If $g\in I$, (T2) places the $Z$-edge of label
$t_1(g)$ at $h$ above layer $5$, hence in layer $6$. If $g\in E$, then
$c_g=5$ and $V_g=\{2,3,4\}$; let $t^\ast=t_2(g)$ and $r=R_{g+b_{t^\ast}}$.
The corner at $h$ with labels $(s^\ast,t^\ast)$ is mapped by $\Phi$ to the
corner at $r$ formed by the $R$-edge of $X_g$ in layer $2$ and the
$Z$-edge at $r$ of label $s^\ast$, both belonging to the pair
$(X_g,Z_{g+a_{s^\ast}+b_{t^\ast}})$. By (T1) the $Z$-edge at $r$ of label
$s_6(g)$ lies below layer $2$, i.e.\ in layer $1$. Since $s_6(g)\neq
s^\ast$ (they label different $L$-edges of $X_g$), the $Z$-edges at $r$
of labels $s^\ast$ and $s_6(g)$ are different edges, and $r$ has exactly
one edge in layer $1$; so the $Z$-edge of label $s^\ast$ at $r$ is not in
layer $1$, hence not below layer $2$, hence the corner at $r$ is
$X$-first. By properness the corner at $h$ is $X$-first: the $Z$-edge at
$h$ of label $t^\ast$ lies above layer $5$, in layer $6$. Distinct sums
enter through (P1): the corner at $h$ has a unique partner, which is the
corner at $r$ just described.

\emph{Step 2: every $R$-qubit with an $X$-edge in layer $5$ has a
$Z$-edge in layer $6$.} Let $r$ carry the layer-$5$ $R$-edge of $X_g$, of
label $t^\ast=t_5(g)$. If $g\in E$ (so $c_g\neq5$), (T1) places the
$Z$-edge of label $s_1(g)$ at $r$ above layer $5$, in layer $6$. If
$g\in I$, then $d_g=5$ and $P_g=\{2,3,4\}$; let $s^\ast=s_2(g)$ and
$h=L_{g+a_{s^\ast}}$. By (T2) the $Z$-edge at $h$ of label $t_6(g)$ is in
layer $1$ and the $L$-edge of $X_g$ occupies layer $2$ at $h$; the
$Z$-edge at $h$ of label $t^\ast\neq t_6(g)$ is a different edge, so it
lies above layer $2$, and the corner at $h$ with labels
$(s^\ast,t^\ast)$ is $X$-first. Its image under $\Phi$ is the corner at
$r=R_{g+b_{t^\ast}}$ formed by the layer-$5$ $R$-edge and the $Z$-edge at
$r$ of label $s^\ast$; by properness it is $X$-first, so that $Z$-edge is
in layer $6$.

\emph{Step 3: cardinalities.} Let $e_5=\#\{g\in E:c_g=5\}$ and
$\iota_5=\#\{g\in I:d_g=5\}$. The $X$-edges in layer $5$ on the $L$
block come from $E$-checks with $c_g=5$ and from $I$-checks with
$d_g\neq5$, so $x_5=e_5+(N-e-\iota_5)$, on as many distinct $L$-qubits.
The $L$-qubits with a $Z$-edge in layer $6$ number $N-x_6=N-e$, since
$x_6=|E|$. Step~1 gives an inclusion of the first set of qubits in the
second, hence
\begin{equation}
e_5+(N-e-\iota_5)\;\le\;N-e .
\end{equation}
The $X$-edges in layer $5$ on the $R$ block come from $E$-checks with
$c_g\neq5$ and $I$-checks with $d_g=5$, so there are $(e-e_5)+\iota_5$ of
them; the $R$-qubits with a $Z$-edge in layer $6$ number $x_6=e$, and
Step~2 gives
\begin{equation}
(e-e_5)+\iota_5\;\le\;e .
\end{equation}
The two inequalities read $e_5\le\iota_5$ and $\iota_5\le e_5$, so
$e_5=\iota_5$ and both inclusions are between sets of equal size, hence
equalities: on the $L$ block a qubit has an $X$-edge in layer $5$ if and
only if it has a $Z$-edge in layer $6$, and likewise on the $R$ block.
At any data qubit each layer holds exactly one edge, $X$ or $Z$. If the
layer-$5$ edge is $X$, the layer-$6$ edge is $Z$; if the layer-$5$ edge
is $Z$, the qubit has no $X$-edge in layer $5$, hence no $Z$-edge in
layer $6$, so its layer-$6$ edge is $X$. Thus the pair $\{5,6\}$ holds one
$X$- and one $Z$-edge at every data qubit.

\emph{Step 4: the other pairs.} Layer reversal $i\mapsto7-i$ maps $\pi$
to a proper schedule with $(a,b)=(1,6)$, maps $E$-checks to $E$-checks
($\{1,6,c\}\mapsto\{1,6,7-c\}$) and $I$-checks to $I$-checks, and
preserves distinct sums; Steps~1--3 applied to the reversed schedule
show that $\{1,2\}$ holds one $X$- and one $Z$-edge at every data qubit
of $\pi$. Each data qubit has three $X$-edges and three $Z$-edges; with
one of each in $\{1,2\}$ and one of each in $\{5,6\}$, the remaining
$X$-edge and $Z$-edge occupy $\{3,4\}$. \hfill$\square$

\paragraph*{Proof of Lemma~\ref{lem:parity}.}
Under the pair structure the $Z$-edges above layer $i$ are one per pair
$\{2j-1,2j\}$ with $2j-1>i$, plus the partner within the pair of $i$ if $i$
is odd, so $\gamma(i)=w-\lceil i/2\rceil+\mathbb{1}[i\text{ odd}]$ and
$\sum_{i=1}^{2w}\gamma(i)=\sum_{j=1}^w(2w-2j+1)=w^2$. \hfill$\square$

\section{Depth six: class schedules, lifting and templates}
\label{si:six}

\paragraph*{Proof of Lemma~\ref{lem:lift}.}
Conflicts: the six edges of $X_g$ receive the values
$T(\varepsilon(g),\cdot,\cdot)$, distinct by (i); the edges at $L_h$ come
from $X_{h-a_s}$ of class $\varepsilon(h)-\alpha_s$ with labels $s$ and from
$Z_{h+b_t}$ of class $\varepsilon(h)+\beta_t$ with labels $t$, distinct by
(ii); $R$-qubits and $Z$-checks likewise. Properness: a pair with
$g'-g=\delta\neq\delta_0$ has one label pair, $Z_{g'}$ has class
$\varepsilon(g)+\alpha_s+\beta_t$, and (iii) gives equal orientations; for
$\delta=\delta_0$ the two collision pairs share the $Z$-check of class
$\varepsilon(g)+\alpha_s+\beta_t=\varepsilon(g)+\alpha_{s'}+\beta_{t'}$
and (iv) gives an even sum. Coincidences among other projected sums play
no role, since every constraint of the code is one of (i)--(iv) for the
labelled pair that defines it. \hfill$\square$

\paragraph*{Normalisation of keys.}
Admissibility of a key is invariant under $(\alpha,\beta)\mapsto
(\alpha+u,\beta+v)$ for any $u,v\in\Zq$: if $T$ is a class schedule for
$(\alpha,\beta)$, then $T'$ defined by $T'(c,\cdot,X\cdot)=T(c,\cdot,X\cdot)$
and $T'(z,\cdot,Z\cdot)=T(z-u-v,\cdot,Z\cdot)$ is a class schedule for
$(\alpha+u,\beta+v)$, because the $L$-qubit of class $c$ in the shifted
system collects $X$-edges from classes $c-\alpha_s-u$ and $Z$-edges from
classes $c+\beta_t+v$, which is the original $L$-qubit of class $c-u$
after relabelling $Z$-classes by $z\mapsto z-u-v$, and the pair
constraints (iii)--(iv) pair $X$-class $c$ with $Z$-class
$c+\alpha_s+\beta_t+u+v$, which the same relabelling sends to the original
pairing; $R$-qubits are handled identically. We therefore fix
$\alpha_1=\beta_1=0$. A designated collision is an unordered pair
$\{(s,t),(s',t')\}$ with $s\neq s'$, $t\neq t'$: three choices of
$\{s,s'\}$, three of $\{t,t'\}$, and two ways to pair them, $18$ in all.
For each, the equality $\alpha_s+\beta_t=\alpha_{s'}+\beta_{t'}$ removes
one of the four free coordinates $(\alpha_2,\alpha_3,\beta_2,\beta_3)$,
leaving $q^3$ keys, hence $18q^3$ keys per modulus, the counts of
Table~\ref{tab:templates-count}. In that enumeration keys related by
label permutations, class translation or layer reversal are counted
separately; the reduction used for $q\ge16$ is described in Methods and
below. The enumerator iterates over the $18$
collisions and the $q^3$ solutions of the equality in lexicographic
order, and its output is checked to contain each normalised key exactly
once. In the proof of Theorem~\ref{thm:six} the key
$(\varepsilon(a_s);\varepsilon(b_t))$ of a code need not be normalised;
the invariance above reduces it to one that is.

\paragraph*{Enumeration.}
For $q\in\{2,3,4,5,6,8\}$, each designated pair and each normalised key,
the constraints (i)--(iv) were encoded on the class system and decided by
SAT; one class schedule was stored per admissible key, checked against
(i)--(iv) and by $\F$ simulation of its pullback to a realising class-C1
code (Methods). $q=7$ was not run; odd $q$ is excluded by
Theorem~\ref{thm:twisted}. Table~\ref{tab:templates-count} and
Fig.~\ref{fig:templates} give the counts. Three regularities are observed
without exception: no admissible key has $\colbar=0$; every admissible key
has $|M|$ even, as Theorem~\ref{thm:twisted} requires, which is why $q=3$
and $q=5$ admit nothing; and for $q\in\{2,4,8\}$ every key with
$\colbar\neq0$ is admissible, which is the first half of
Proposition~\ref{prop:template}. The refinement that a key is admissible
exactly when $\operatorname{ord}(\colbar)$ is even is false: at $q=6$,
$468$ admissible keys have $\operatorname{ord}(\colbar)=3$.

\paragraph*{Keys with $\colbar=q/2$ for $q\ge16$.}
By the reduction in Methods the keys with $\colbar=q/2$ are
$\kappa_q(a,b)=(0,q/2,a;0,q/2,b)$, $(a,b)\in\Zq^2$, modulo
$(a,b)\sim(b,a)\sim(\omega a,\omega b)\sim(a+q/2,b+q/2)$, and every such
key reduces to a member of the family $\kappa_{q'}(1,b)$ at a modulus
$q'\mid q$. Table~\ref{tab:family} gives the numbers of representatives
and the outcome of the SAT decisions. For $q\le32$ all representatives
were decided without assumptions ($39$ at $q=32$ in $17$ minutes); single
instances at $q=64$ take up to half an hour, and the family was decided
there and beyond under the half-period symmetry described next, which the
solver handles in seconds.

\paragraph*{A universal half-period symmetry.}
Imposing $T(c+q/2,e)=\sigma(T(c,e))$ for a single solver-chosen
permutation $\sigma$ of the six layers is satisfiable only for $b=5$ at
$q=8$ and $b=9$ at $q=16$. Imposing it with a permutation
$\sigma_{\mathrm{type}(e)}$ depending on the edge type ($XL$, $XR$, $ZL$,
$ZR$) is satisfiable for every $b$ at $q=8,16,32,64,128$, with
fixed-point-free involutions in every solution inspected. One assignment
serves all $b$ at once: a
joint SAT instance over all $b$ at $q=16$ with the permutation depending
only on the block gives $\sigma_L=(13)(24)(56)$ on both $L$-edge types and
$\sigma_R=(12)(36)(45)$ on both $R$-edge types, and with these two
involutions imposed every member of the family at $q=2,4,\dots,512$ admits
a class schedule; this is the computation behind
Proposition~\ref{prop:template} for $q\ge64$ (all $512$ keys at $q=512$ in
twelve minutes). A second pair, $\sigma_L=(14)(23)(56)$,
$\sigma_R=(12)(35)(46)$, serves all $b$ for $q\le16$; a permutation shared
by $L$ and $R$, or one shared by $XL$ and $ZR$ and another by $XR$ and
$ZL$, does not exist at $q=8$ or $q=16$. The symmetry cannot be refined:
requiring in addition $T(c+q/4,e)=\sigma'_{\mathrm{type}(e)}(T(c,e))$ is
unsatisfiable for every $b$ at $q=8$ and $q=16$, and a schedule depending
on $c$ only through $(c\bmod2,\lfloor4c/q\rfloor)$ exists for $4$ of the
$16$ values of $b$ at $q=16$, one depending only on
$(c\bmod4,\mathbb{1}[c\ge q/2])$ for $2$ of them. An analytic template
therefore has to prescribe the layers on the classes $c<q/2$; the
half-period symmetry then determines the rest.



\paragraph*{Proof of Lemma~\ref{lem:height}.}
Put $q=2^{h+1}$ and $\bar G=G/qG\cong\bigoplus_i\mathbb{Z}_{2^{e_i}}$ with
$e_i\le h+1$. The image $\bar{\col}$ of $\col$ in $\bar G$ is non-zero
because $\col\notin qG$, and $2\bar{\col}=0$ because $\col\in2^hG$; hence
$\bar{\col}=\sum_ic_i$ with $c_i\in\{0,2^{e_i-1}\}$ and some
$c_{i_0}\neq0$. Composing $G\to\bar G\to\mathbb{Z}_{2^{e_{i_0}}}$ with the
injection $x\mapsto2^{h+1-e_{i_0}}x$ into $\Zq$ sends $\col$ to
$2^{h+1-e_{i_0}}\cdot2^{e_{i_0}-1}=2^h$. \hfill$\square$

\begin{table}[!htbp]
\caption{\textbf{Class schedules by modulus.} Normalised keys with one
designated collision, keys admitting a class schedule, and $|M|$ among
them; $q=7$ was not enumerated (odd, excluded by
Theorem~\ref{thm:twisted}).}
\label{tab:templates-count}
\begin{ruledtabular}
\begin{tabular}{cccl}
$q$ & keys & admissible & $|M|$ among admissible \\
\colrule
2 & 144  & 72   & $\{2\}$ \\
3 & 486  & 0    & -- \\
4 & 1152 & 864  & $\{2,4\}$ \\
5 & 2250 & 0    & -- \\
6 & 3888 & 2412 & $\{2,6\}$ \\
8 & 9216 & 8064 & $\{2,4,8\}$ \\
\end{tabular}
\end{ruledtabular}
\end{table}

\begin{figure}[!htbp]
\centering
\includegraphics[width=\SIcolumnfigurewidth]{figS3_templates_by_q.pdf}
\caption{\textbf{Class schedules by modulus.} For $q\in\{2,3,4,5,6,8\}$,
stacked counts of keys that admit a class schedule, do not admit one with
$\colbar\neq0$, and have $\colbar=0$.}
\label{fig:templates}
\end{figure}

\begin{table}[!htbp]
\caption{\textbf{Keys with $\colbar=q/2$.} Number of keys
$\kappa_q(a,b)$, number of representatives under the three symmetries, and
how the keys were decided: all representatives by unconstrained SAT
($q\le32$), or every member of the family~\eqref{eq:family} under the
half-period symmetry ($q\ge64$; also done for $q\le32$).}
\label{tab:family}
\begin{ruledtabular}
\begin{tabular}{cccl}
$q$ & keys & reps. & decided \\
\colrule
2   & 4     & 2   & all, SAT \\
4   & 16    & 5   & all, SAT \\
8   & 64    & 11  & all, SAT \\
16  & 256   & 21  & all, SAT \\
32  & 1024  & 39  & all, SAT \\
64  & 4096  & 73  & family with symmetry, SAT \\
128 & 16384 & 139 & family with symmetry, SAT \\
256 & 65536 & 269 & family with symmetry, SAT \\
512 & 262144 & 527 & family with symmetry, SAT \\
\end{tabular}
\end{ruledtabular}
\end{table}

\paragraph*{The $q=2$ template.}
The template displayed in Fig.~\ref{fig:mechanism} is the admissible
$q=2$ key $(0,1,0;0,1,0)$ with collision $\{(1,1),(2,2)\}$; its $24$ layer
values are listed in Table~\ref{tab:template-m2}. Pulled back along
$\varepsilon(g)=g\bmod2$ on $G=\mathbb{Z}_{14}$ with $A=\{0,1,10\}$,
$B=\{0,13,2\}$, it is a proper depth-$6$ schedule of a connected class-C1
code with parameters $[[28,6,4]]$ (the distance is exact: no logical
operator of weight below $4$ exists, by exhaustive search), verified by
$\F$ simulation.
Fifteen further pullbacks, across $q\in\{2,4,6,8\}$ and the groups
$\mathbb{Z}_{10}$, $\mathbb{Z}_{12}$, $\mathbb{Z}_{16}$, $\mathbb{Z}_{18}$,
$\mathbb{Z}_{24}$, $\mathbb{Z}_{32}$, $\mathbb{Z}_2\times\mathbb{Z}_6$ and
$\mathbb{Z}_2\times\mathbb{Z}_8$, with projections merging three to six
further sums, all pass.

\begin{table}[!htbp]
\caption{\textbf{The $q=2$ class schedule of Fig.~\ref{fig:mechanism}.}
Layers by class $c$, edge type and label, for the key $(0,1,0;0,1,0)$ with
collision $\{(1,1),(2,2)\}$; pulled back along $\varepsilon(g)=g\bmod2$ on
$\mathbb{Z}_{14}$ with $A=\{0,1,10\}$, $B=\{0,13,2\}$, it is a proper
depth-$6$ schedule of a $[[28,6,4]]$ code, verified by $\F$
simulation.}
\label{tab:template-m2}
\begin{ruledtabular}
\begin{tabular}{llccc}
class & edge type & label 1 & label 2 & label 3 \\
\colrule
$c=0$ & $X$-check, $L$-edges & 3 & 6 & 1 \\
      & $X$-check, $R$-edges & 4 & 2 & 5 \\
      & $Z$-check, $L$-edges & 5 & 3 & 2 \\
      & $Z$-check, $R$-edges & 1 & 4 & 6 \\
\colrule
$c=1$ & $X$-check, $L$-edges & 5 & 4 & 2 \\
      & $X$-check, $R$-edges & 1 & 3 & 6 \\
      & $Z$-check, $L$-edges & 4 & 6 & 1 \\
      & $Z$-check, $R$-edges & 3 & 2 & 5 \\
\end{tabular}
\end{ruledtabular}
\end{table}

\paragraph*{The $2$-height in the data set.}
A homomorphism $G\to\mathbb{Z}_{2^j}$ nonzero on $\col$ exists iff
$\col\notin2^jG$, since $\mathrm{Hom}(G/2^jG,\mathbb{Z}_{2^j})$ separates
the points of $G/2^jG$; restricting Theorem~\ref{thm:six} to $q\le8$ would
therefore miss even-order elements inside $8G$, such as $\col=8$ in
$\mathbb{Z}_{32}$, which has $2$-height $3$ and is covered through $q=16$.
Table~\ref{tab:twoheight} splits the $598$ class-C1 instances by the
$2$-height; the two instances with $h(\col)=3$, both on $\mathbb{Z}_{32}$,
were also found to admit $\mathbb{Z}_{16}$-twisted proper schedules by
direct search, verified by simulation.

\begin{table}[!htbp]
\caption{\textbf{Class-C1 instances by the $2$-height of $\col$.} The two
instances with $h(\col)=3$ also admit $\mathbb{Z}_{16}$-twisted witnesses
found by direct search.}
\label{tab:twoheight}
\begin{ruledtabular}
\begin{tabular}{lccl}
condition on $\col$ & instances & witness & verdict \\
\colrule
$h(\col)\le2$ ($\col\notin8G$) & 263 & template, $q\le8$ & all SAT \\
$h(\col)=3$ ($\col\in8G$) & 2 & template, $q=16$ & both SAT \\
$\operatorname{ord}(\col)$ odd & 333 & -- & Table~\ref{tab:c1} \\
\end{tabular}
\end{ruledtabular}
\end{table}

\section{Sharpness data}
\label{si:sharpness}

\paragraph*{Overlap taxonomy.}
Let $n_4$ and $n_6$ be the numbers of check pairs sharing exactly four and
six qubits; then $n_4+3n_6=N\sum_\delta\binom{S(\delta)}{2}$, and D is
$(0,0)$, C1 is $(N,0)$, C2 is $(2N,0)$ or $(0,N)$, any other pattern
``other''. The instance $G=\mathbb{Z}_8$, $A=\{0,1,4\}$, $B=\{0,7,4\}$ has
multiplicities $(1,1,1,1,2,3)$, hence $n_4=n_6=N$, and is ``other''
although exactly one $\delta$ has $S(\delta)=2$; this is why C1 requires
$S\le1$ elsewhere. The classifier reads the label from $H_XH_Z^{\mathsf T}$
and agrees with the sums on every classified instance and on this example.

\paragraph*{Small groups.}
The sweep over even orders $2$--$8$ is exhaustive over pairs with $0\in A$,
$0\in B$. Under $(A,B)\sim(\phi(A)+t_1,\phi(B)+t_2)$ for
$\phi\in\mathrm{Aut}(G)$ and $(A,B)\sim(B,A)$, the $1441$ pairs fall into
$33$ orbits, $32$ satisfiable and one unsatisfiable, the latter being the
$\mathbb{Z}_6$ instance $A=B=\{0,2,4\}$ alone (Tables~\ref{tab:sweep},
\ref{tab:sweep-census}); no orbit contains both outcomes. That instance has
components of six data qubits with $\sum_vr_vs_v=54\equiv2\pmod4$, so
Theorem~\ref{thm:mod4} excludes it, and its DRAT certificate confirms the
theorem. The odd orders $3$, $5$, $7$, $9$ were swept up to the same
equivalence: $1830$ pairs in $27$ orbits, decided under a $20$-second cap
(Table~\ref{tab:sweep-odd}). The three distinct-sums orbits, all of order
$9$ (recorded as O9-a, O9-b, O9-c), are unsatisfiable; of the $24$ orbits
with coinciding sums, $2$ are excluded by Theorem~\ref{thm:mod4} and $22$
reached the cap, so this part of the sweep is inconclusive rather than
negative.

\begin{table}[!htbp]
\caption{\textbf{Exhaustive sweep over even $|G|\le8$.}}
\label{tab:sweep}
\begin{ruledtabular}
\begin{tabular}{lc}
 & count \\
\colrule
normalised instances & 1441 \\
proper depth-$6$ schedule found and verified & 1440 \\
unsatisfiable, excluded by Theorem~\ref{thm:mod4} & 1 \\
\end{tabular}
\end{ruledtabular}
\end{table}

\begin{table}[!htbp]
\caption{\textbf{Census of the even-order sweep.} Orbits under
automorphisms, translations and block swap, and outcomes of the
representatives.}
\label{tab:sweep-census}
\begin{ruledtabular}
\begin{tabular}{lrrrrr}
group & $|\mathrm{Aut}(G)|$ & pairs & orbits & SAT & UNSAT \\
\colrule
$\mathbb{Z}_4$            & 2   & 9    & 1  & 1  & 0 \\
$\mathbb{Z}_6$            & 2   & 100  & 7  & 6  & 1 \\
$\mathbb{Z}_8$            & 4   & 441  & 13 & 13 & 0 \\
$\mathbb{Z}_2\times\mathbb{Z}_2$ & 6 & 9  & 1  & 1  & 0 \\
$\mathbb{Z}_2\times\mathbb{Z}_4$ & 8 & 441 & 9 & 9  & 0 \\
$\mathbb{Z}_2^3$          & 168 & 441  & 2  & 2  & 0 \\
\colrule
total                     &     & 1441 & 33 & 32 & 1 \\
\end{tabular}
\end{ruledtabular}
\end{table}

\begin{table}[!htbp]
\caption{\textbf{Odd-order sweep.} Orbits under the same equivalence and
outcomes of the representatives at a $20$-second cap.}
\label{tab:sweep-odd}
\begin{ruledtabular}
\begin{tabular}{lrrrrr}
group & orbits & pairs & SAT & UNSAT & capped \\
\colrule
$\mathbb{Z}_3$ & 1  & 1   & 0 & 0 & 1  \\
$\mathbb{Z}_5$ & 2  & 36  & 0 & 0 & 2  \\
$\mathbb{Z}_7$ & 5  & 225 & 0 & 1 & 4  \\
$\mathbb{Z}_9$ & 12 & 784 & 0 & 2 & 10 \\
$\mathbb{Z}_3\times\mathbb{Z}_3$ & 7 & 784 & 0 & 2 & 5 \\
\colrule
total & 27 & 1830 & 0 & 5 & 22 \\
\end{tabular}
\end{ruledtabular}
\end{table}

\begin{figure}[!htbp]
\centering
\includegraphics[width=\SIcolumnfigurewidth]{figS2_sweep_by_group.pdf}
\caption{\textbf{Depth-$6$ outcomes on groups of order at most nine.}
Per-group counts of normalised instances by outcome; odd orders are swept
up to equivalence.}
\label{fig:sweep}
\end{figure}

\paragraph*{Instance generation for the sharpness map.}
The sampled instances of Table~\ref{tab:h2} are not enumerated. For cyclic
groups of order $N\in\{10,\dots,24,28,30,32,36\}$, $A$ and $B$ were each
drawn as $\{0\}$ plus two elements uniformly from $\{1,\dots,N-1\}$ (seed
$7$), classified, and kept up to a fixed quota per $(N,\text{class})$; the
subgroup instances $A=B=\{0,N/3,2N/3\}$ were added for $N\in\{12,18\}$.
Duplicates were removed by literal equality only. Under the equivalence
above, the $664$ instances fall into $475$ orbits ($416$ of $598$ in C1,
$30$ of $35$ in D, $26$ of $28$ in C2, all $3$ subgroup-type), so counts
over the $664$ are weighted by orbit size, most heavily at small $N$;
per-group counts are in the repository.

\paragraph*{Sampled groups and the seven codes of Ref.~\cite{Bravyi2024}.}
Table~\ref{tab:h2} lists the $664$ instances by class; the $\mathbb{Z}_6$
instance is the only one of order below $10$. Every distinct-sums instance
is unsatisfiable, as Theorem~\ref{thm:main} requires; the smallest and
largest, at $|G|=10$ and $24$, and the seven codes of
Ref.~\cite{Bravyi2024} carry DRAT certificates. Of the class-C2 instances
$27$ admit depth $6$ and one is excluded by Theorem~\ref{thm:mod4}; both
subgroup instances, in which $A\cup B$ generates a proper subgroup, are
excluded by Theorem~\ref{thm:mod4} applied to their components.

\begin{table}[!htbp]
\caption{\textbf{The $664$ depth-$6$ instances by class.} S: proper
depth-$6$ schedule found and verified; U$_{\rm t}$: excluded by
Theorem~\ref{thm:mod4} or~\ref{thm:main}; U$_{\rm s}$: refuted by the
solver (of the $62$, six carry DRAT certificates and $56$ have no proof
logging); O: solver capped, no theorem applies.}
\label{tab:h2}
\begin{ruledtabular}
\begin{tabular}{lrrrrr}
class & $n$ & S & U$_{\rm t}$ & U$_{\rm s}$ & O \\
\colrule
D ($|G|=10$--$24$)  & 28  & 0   & 28  & 0  & 0  \\
D (Ref.~\cite{Bravyi2024}) & 7 & 0 & 7  & 0  & 0  \\
C1                  & 598 & 295 & 211 & 62 & 30 \\
C2                  & 28  & 27  & 1   & 0  & 0  \\
subgroup type       & 2   & 0   & 2   & 0  & 0  \\
$\mathbb{Z}_6$, $A=B=\{0,2,4\}$ & 1 & 0 & 1 & 0 & 0 \\
\colrule
total               & 664 & 322 & 250 & 62 & 30 \\
\end{tabular}
\end{ruledtabular}
\end{table}

\paragraph*{Class C1.}
Table~\ref{tab:c1} gives the $598$ class-C1 instances by the parity of
$\operatorname{ord}(\col)$, the $2$-height of $\col$,
Theorem~\ref{thm:mod4}, and outcome. All $265$ even-order instances admit
depth $6$ by Theorem~\ref{thm:six} ($h(\col)\le3$ throughout): $262$ were
also found by the solver, one by a template witness beyond the cap, and
the two with $h(\col)=3$ by $\mathbb{Z}_{16}$-twisted witnesses as well as
by the $q=16$ templates. Of the $333$ odd-order instances $211$ are
excluded by Theorem~\ref{thm:mod4}, the solver confirming $131$ and
reaching its cap on $80$; of the remaining $122$, $30$ admit depth $6$,
$62$ are refuted by the solver ($6$ with DRAT certificates, $56$ without
proof logging), and $30$ are open. Fig.~\ref{fig:periods} shows the
smallest modulus $q^\ast$ under which a found proper schedule is twisted:
$24$ of the $30$ satisfiable odd-order instances have an even
$q^\ast\in\{6,10,12,14\}$, and for the other $6$ the found schedule has no
proper period and the search for a twisted schedule over every $q\mid N$
completed without a cap and found none, so those codes admit no periodic
proper schedule. Fig.~\ref{fig:thm1profile} shows how
Theorem~\ref{thm:mod4} distributes over the sweep.

\begin{table}[!htbp]
\caption{\textbf{Class-C1 instances.} obs.: Theorem~\ref{thm:mod4}
obstructs; S$_{\rm s}$ satisfiable by solver; S$_{\rm w}$ by twisted
witness; U unsatisfiable (in the last row, $6$ with DRAT certificates and
$56$ solver-reported); O open. In the obstructed row the solver confirmed
$131$ and reached its cap on $80$.}
\label{tab:c1}
\begin{ruledtabular}
\begin{tabular}{lcrrrrr}
$\col$ & obs. & $n$ & S$_{\rm s}$ & S$_{\rm w}$ & U & O \\
\colrule
even, $h(\col)\le2$ & no  & 263 & 262 & 1 & 0   & 0  \\
even, $h(\col)=3$   & no  & 2   & 0   & 2 & 0   & 0  \\
odd              & yes & 211 & 0   & 0 & 211 & 0  \\
odd              & no  & 122 & 25  & 5 & 62  & 30 \\
\colrule
total            &     & 598 & 287 & 8 & 273 & 30 \\
\end{tabular}
\end{ruledtabular}
\end{table}

\begin{figure}[!htbp]
\centering
\includegraphics[width=\SIcolumnfigurewidth]{figS4_periods.pdf}
\caption{\textbf{Periods of satisfiable class-C1 schedules.} Smallest
modulus $q^\ast$ under which a found proper schedule is twisted, split by
the parity of $\operatorname{ord}(\col)$; the last category collects the
six instances whose found schedule has no proper period and for which no
twisted schedule exists at any $q\mid N$.}
\label{fig:periods}
\end{figure}

\begin{figure}[!htbp]
\centering
\includegraphics[width=\SIcolumnfigurewidth]{figS5_theorem1_profile.pdf}
\caption{\textbf{Theorem~\ref{thm:mod4} across the sweep.} Instances by
$\sum_vr_vs_v\bmod4$ per component, which for $r_v=s_v=3$ equals
$9n/2\bmod4$ with $n$ the component size, and by depth-$6$ outcome.}
\label{fig:thm1profile}
\end{figure}

\paragraph*{Coverage.}
Table~\ref{tab:coverage} partitions the $664$ instances by mathematical
and solver status. Mathematical status takes precedence: an instance
excluded by Theorem~\ref{thm:mod4} or~\ref{thm:main} counts as
unsatisfiable whether or not the solver finished, and ten of them also
carry DRAT certificates; five further instances not covered by a theorem
are certified; the eight capped satisfiable instances, three with
$\operatorname{ord}(\col)$ even and five odd, are settled by twisted
witnesses.

\begin{table}[!htbp]
\caption{\textbf{Coverage of the $664$ instances.} Rows: mathematical
status; columns: solver status.}
\label{tab:coverage}
\begin{ruledtabular}
\begin{tabular}{lcccc}
status & unsat & sat & capped & total \\
\colrule
unsatisfiable by theorem      & 170 & 0   & 80  & 250 \\
unsatisfiable by certificate  & 6   & 0   & 0   & 6   \\
unsatisfiable by solver only  & 56  & 0   & 0   & 56  \\
open                          & 0   & 0   & 30  & 30  \\
satisfiable (verified)        & 0   & 314 & 8   & 322 \\
\colrule
total                         & 232 & 314 & 118 & 664 \\
\end{tabular}
\end{ruledtabular}
\end{table}

\section{Certificates and validation}
\label{si:certs}

\paragraph*{Encoding.}
Variables $x_{e,i}$ per Tanner edge $e$ and layer $i\in[D]$; exactly-one
clauses per edge; at-most-one clauses for every Tanner vertex (data qubits
and check ancillas) and layer; for each overlapping pair,
Eq.~\eqref{eq:parity} with the orientation of a corner $(q,e,e')$
represented by a Tseitin variable for
$\bigvee_{i<j}(x_{e,i}\wedge x_{e',j})$ and the parity by Tseitin variables
for the partial sums; $10^3$--$10^4$ variables for the instances here.
Class schedules use the same clauses on the class system with (i)--(iv)
in place of the code's constraints. The template encoder imposes
(iii)--(iv) per labelled pair, never regrouped by projected sum. An early
version added a symmetry-breaking clause that fixes the orientation of one
designated corner $(q_0,e_0,e_0')$ to $X$-first, i.e.\
$\bigvee_{i<j}(x_{e_0,i}\wedge x_{e_0',j})$. Layer reversal maps proper
schedules to proper schedules and flips every orientation, so every proper
schedule or its reversal satisfies the clause; the augmented formula is
therefore unsatisfiable if and only if the original is, and a certificate
for the augmented formula proves unsatisfiability of the original.
Certificates for the seven codes of Ref.~\cite{Bravyi2024} were
regenerated without the clause; the $\mathbb{Z}_6$ certificate and the two
class-D certificates retain it.

\paragraph*{Solvers and certificates.}
\textsc{kissat}~4.0.4~\cite{Biere2024Kissat} with DRAT logging, proofs
verified by \textsc{drat-trim}~\cite{Wetzler2014DRAT,HeuleHuntWetzler2013};
\textsc{CaDiCaL}~\cite{Biere2024CaDiCaL} and \textsc{Glucose}~4.2 through
PySAT~\cite{Ignatiev2018PySAT,AudemardSimon2009}; every satisfying
assignment decoded and verified by $\F$ simulation; every capped run
recorded as capped. Table~\ref{tab:certs-full} lists all eighteen certified
instances; Fig.~\ref{fig:certs} shows solve times and proof sizes. All
runs were single-threaded on one workstation (Intel Core Ultra 9 275HX,
$24$ logical cores, $15.6$~GiB, Linux/WSL2); wall and CPU time agree
within a few per cent.

\begin{table*}[!t]
\caption{\textbf{DRAT certificates.} Instance, whether the
symmetry-breaking clause was present, \textsc{kissat} solve time, and
proof size; all certificates were verified by \textsc{drat-trim}.}
\label{tab:certs-full}
\centering

\begin{ruledtabular}
\begin{tabular}{lccc}
instance & clause & time (s) & proof \\
\colrule
$[[72,12,6]]$      & no  & 4.0 & 24 MB \\
$[[90,8,10]]$      & no  & 4.2 & 41 MB \\
$[[108,8,10]]$     & no  & 4.8 & 52 MB \\
$[[144,12,12]]$    & no  & 3.6 & 35 MB \\
$[[288,12,18]]$    & no  & 1.4 & 18 MB \\
$[[360,12,\le24]]$ & no  & 6.3 & 68 MB \\
$[[756,16,\le34]]$ & no  & 3.3 & 46 MB \\
D, $\mathbb{Z}_{10}$, $A=\{0,3,6\}$, $B=\{0,1,9\}$
  & yes & 4.1 & 13 MB \\
D, $\mathbb{Z}_{24}$, $A=\{0,1,8\}$, $B=\{0,5,14\}$
  & yes & 3.6 & 18 MB \\
$\mathbb{Z}_6$, $A=B=\{0,2,4\}$
  & yes & 558 & 2.7 GB \\
6 class-C1, odd order, not obstructed
  & yes & 16--163 & 123--911 MB \\
C1 odd, $\mathbb{Z}_{45}$, $A=\{0,12,33\}$, $B=\{0,24,39\}$
  & yes & 967 & 2.4 GB \\
C1 odd, $\mathbb{Z}_{15}\times\mathbb{Z}_3$, $A=\{(0,0),(8,2),(13,1)\}$,
$B=\{(0,0),(1,1),(14,2)\}$
  & yes & 600 & 2.8 GB \\
\end{tabular}
\end{ruledtabular}
\end{table*}

\begin{figure*}[!t]
\centering
\includegraphics[width=\SIwidefigurewidth]{figS6_certificates.pdf}
\caption{\textbf{Certificates.} Solve time and DRAT proof size for each
certified instance, filled markers with the layer-reversal clause and open
markers without; Table~\ref{tab:certificates} reports the values without
the clause.}
\label{fig:certs}
\end{figure*}

\paragraph*{Cross-check.}
The $17$ certificate CNFs of the first ten instances (seven codes of
Ref.~\cite{Bravyi2024} in both encodings, the two class-D instances, the
$\mathbb{Z}_6$ instance) were re-decided by \textsc{CaDiCaL} and
\textsc{Glucose} with a $300$-second cap: $16$ are unsatisfiable under
both; the $\mathbb{Z}_6$ file is decided by \textsc{CaDiCaL} in $233$
seconds and not by \textsc{Glucose} within the cap.

\paragraph*{Model ablation.}
The parity criterion~\eqref{eq:parity} and direct $\F$ simulation agree on
$300$ random valid depth-$\Delta$ layerings of each of $[[72,12,6]]$ and
\texttt{oe23}, $600$ in all; the identity
$\sum_vN_v=\frac12\sum_vr_vs_v$ of Theorem~\ref{thm:mod4} holds on $200$
random layerings of each, $400$ in all, proper or not.

\paragraph*{The uniform depth-$7$ schedule.}
The schedule of Proposition~\ref{prop:seven} is Table~\ref{tab:depth7main}
of the main text. Its nine corner orientations are
$o_L=o_R=\mathbb{1}[s\le2]$; it was simulated on the seven codes of
Ref.~\cite{Bravyi2024}, $20$ random distinct-sums instances and $20$
random class-C1/C2 instances, $47$ of $47$ proper. It is not equivalent to
the layer order of Ref.~\cite{Bravyi2024} under label permutations, layer
reversal, $X$/$Z$ duality and block swap (Methods): the published cycle has
$o_L=o_R=\mathbb{1}[s=2]$, and the $8496$ translation-invariant proper
depth-$7$ schedules fall into $41$ orbits under that group, the two
schedules lying in different orbits.

\paragraph*{Lemma checks.}
Four codes with one even weight admit proper depth-$\Delta$ schedules
(Table~\ref{tab:lemmas}); all such schedules were enumerated up to $200$
per instance, $417$ in all ($5$ translation-invariant, $72$ twisted, $340$
without a proper period). On every schedule: $a+b=\Delta+1$; every
$X$-check with an $L$-edge in layer $a$ has one in layer $b$ and
conversely; $\Phi$ preserves orientation; $F_L=F_R$; (T1) and (T2) in
shell form; and $F_L+F_R$ even. The conditional statement
$F_L+F_R=\wA\wB$ when no layer is mixed applies to $140$ of the $417$ and
holds on all of them. \texttt{oe12} and \texttt{oe21} are not dual: the
dual of \texttt{oe12} has $A=\{0,4\}$, not equivalent to $\{0,3\}$ under
the units of $\mathbb{Z}_8$.

\begin{table}[!htbp]
\caption{\textbf{Lemma checks.}}
\label{tab:lemmas}
\begin{ruledtabular}
\begin{tabular}{lcccccc}
instance & $G$ & $A$ & $B$ & weights & $\Delta$ & schedules \\
\colrule
\texttt{oe12} & $\mathbb{Z}_8$    & $\{0\}$       & $\{0,4\}$   & $(1,2)$ & 3 & 16  \\
\texttt{oe21} & $\mathbb{Z}_8$    & $\{0,3\}$     & $\{0\}$     & $(2,1)$ & 3 & 1   \\
\texttt{oe23} & $\mathbb{Z}_{12}$ & $\{0,1\}$     & $\{0,4,8\}$ & $(2,3)$ & 5 & 200 \\
\texttt{oe43} & $\mathbb{Z}_{16}$ & $\{0,1,2,3\}$ & $\{0,4,8\}$ & $(4,3)$ & 7 & 200 \\
\colrule
total & & & & & & 417 \\
\end{tabular}
\end{ruledtabular}
\end{table}

\paragraph*{Coset identity.}
$\sum_{c\in C}(F_L(c)+F_R(c))=|M|\wA\wB$ was evaluated on every twisted
object in the repository: $12{,}931$ cosets, $295$ from the $289$ twisted
witnesses stored with their homomorphism and modulus, and $12{,}636$ from
the $11{,}412$ admissible template keys ($10{,}404$ with a single coset,
$1{,}008$ with several), with no failure and $|M|$ even throughout. On the
$27$ satisfiable class-C2 instances, twisted schedules were enumerated for
every $q\mid N$ under a $20$-second cap: $323$ found, $|M|$ taking eleven
even values between $2$ and $24$, $29$ combinations capped. Both subgroup
instances are unsatisfiable, so beyond classes C1 and C2 the theorem is
supported by its proof only.

\paragraph*{Reproducibility.}
The repository contains the encoder, the instance generator, all
certificates, schedules and witnesses with their homomorphisms, the
templates, the simulation checker, one CSV per table and figure, a
manifest mapping every figure to its columns, and a script that checks
every number of Sections~\ref{sec:six}--\ref{sec:methods} against those
files.

\section{Depth-six codes: search data and explicit schedules}
\label{si:codes}

\paragraph*{Search data.}
Tables~\ref{tab:codes-a}--\ref{tab:codes-c} list, for every group and
every class (D: distinct sums; C1, even: class C1 with
$\operatorname{ord}(\col)$ even; C1, odd: class C1 with
$\operatorname{ord}(\col)$ odd), the Pareto-optimal $(k,d)$ pairs with
$d\ge4$ found by the search of Methods, one representative $(A,B)$ per
pair (elements of $\mathbb{Z}_\ell\times\mathbb{Z}_m$ written $(i,j)$, of
$\mathbb{Z}_N$ written $i$), the collision element, its order and its
$2$-height. A distance without ``$\le$'' is exact (meet-in-the-middle
confirmation); the search is exhaustive over all pairs with $0\in A$,
$0\in B$ for $N\le36$ and samples $4\times10^5$ pairs per group for
$N\in\{45,54,72\}$ (about $21\%$ of the pairs at $N=54$ and $6.5\%$ at
$N=72$), so the fronts at $N\ge45$ are lower bounds on what an exhaustive
search would find. Every class-C1 code of even order in the tables has
$h(\col)\le2$ and is covered by Theorem~\ref{thm:six}; the groups
$\mathbb{Z}_6\times\mathbb{Z}_6$, $\mathbb{Z}_2\times\mathbb{Z}_{18}$,
$\mathbb{Z}_9\times\mathbb{Z}_6$, $\mathbb{Z}_{18}\times\mathbb{Z}_3$ and
$\mathbb{Z}_{54}$ have $2G=4G$, so every even-order element there has
$2$-height $0$.

\begin{table*}[!t]
\caption{\textbf{Pareto fronts, $N\in\{18,24,27\}$.}}
\label{tab:codes-a}
\footnotesize
\begin{ruledtabular}
\begin{tabular}{llllllcc}
group & class & $[[n,k,d]]$ & $A$ & $B$ & $\col$ & $\operatorname{ord}(\col)$ & $h(\col)$ \\
\colrule
$\mathbb{Z}_{18}$ & C1, even & $[[36,4,6]]$ & $\{0,10,17\}$ & $\{0,1,14\}$ & 17 & 18 & 0 \\
$\mathbb{Z}_{18}$ & C1, odd & $[[36,4,6]]$ & $\{0,1,2\}$ & $\{0,5,16\}$ & 2 & 9 & -- \\
$\mathbb{Z}_{18}$ & D & $[[36,4,6]]$ & $\{0,8,13\}$ & $\{0,2,4\}$ & -- & -- & -- \\
$\mathbb{Z}_{2}\times\mathbb{Z}_{12}$ & C1, even & $[[48,8,4]]$ & $\{(0,0),(0,1),(1,2)\}$ & $\{(0,0),(0,5),(1,7)\}$ & (1,2) & 6 & 0 \\
$\mathbb{Z}_{2}\times\mathbb{Z}_{12}$ & C1, even & $[[48,4,6]]$ & $\{(0,0),(0,7),(1,2)\}$ & $\{(0,0),(0,7),(0,11)\}$ & (0,7) & 12 & 0 \\
$\mathbb{Z}_{2}\times\mathbb{Z}_{12}$ & C1, odd & $[[48,8,4]]$ & $\{(0,0),(0,8),(1,7)\}$ & $\{(0,0),(0,2),(0,10)\}$ & (0,8) & 3 & -- \\
$\mathbb{Z}_{2}\times\mathbb{Z}_{12}$ & D & $[[48,8,4]]$ & $\{(0,0),(1,1),(1,8)\}$ & $\{(0,0),(0,8),(0,10)\}$ & -- & -- & -- \\
$\mathbb{Z}_{2}\times\mathbb{Z}_{12}$ & D & $[[48,4,6]]$ & $\{(0,0),(0,1),(0,5)\}$ & $\{(0,0),(1,8),(1,10)\}$ & -- & -- & -- \\
$\mathbb{Z}_{4}\times\mathbb{Z}_{6}$ & C1, even & $[[48,8,4]]$ & $\{(0,0),(1,4),(2,5)\}$ & $\{(0,0),(0,5),(3,4)\}$ & (1,1) & 12 & 0 \\
$\mathbb{Z}_{4}\times\mathbb{Z}_{6}$ & C1, even & $[[48,4,6]]$ & $\{(0,0),(0,1),(2,2)\}$ & $\{(0,0),(1,4),(2,2)\}$ & (2,2) & 6 & 1 \\
$\mathbb{Z}_{4}\times\mathbb{Z}_{6}$ & C1, odd & $[[48,8,4]]$ & $\{(0,0),(0,1),(0,2)\}$ & $\{(0,0),(0,4),(2,5)\}$ & (0,2) & 3 & -- \\
$\mathbb{Z}_{4}\times\mathbb{Z}_{6}$ & D & $[[48,8,4]]$ & $\{(0,0),(2,5),(3,4)\}$ & $\{(0,0),(0,4),(2,2)\}$ & -- & -- & -- \\
$\mathbb{Z}_{4}\times\mathbb{Z}_{6}$ & D & $[[48,4,6]]$ & $\{(0,0),(2,1),(3,5)\}$ & $\{(0,0),(1,2),(2,4)\}$ & -- & -- & -- \\
$\mathbb{Z}_{24}$ & C1, even & $[[48,8,4]]$ & $\{0,10,20\}$ & $\{0,20,22\}$ & 20 & 6 & 2 \\
$\mathbb{Z}_{24}$ & C1, even & $[[48,4,8]]$ & $\{0,10,11\}$ & $\{0,2,16\}$ & 10 & 12 & 1 \\
$\mathbb{Z}_{24}$ & C1, odd & $[[48,8,4]]$ & $\{0,2,16\}$ & $\{0,4,8\}$ & 16 & 3 & -- \\
$\mathbb{Z}_{24}$ & C1, odd & $[[48,4,6]]$ & $\{0,1,8\}$ & $\{0,8,19\}$ & 8 & 3 & -- \\
$\mathbb{Z}_{24}$ & D & $[[48,4,8]]$ & $\{0,10,23\}$ & $\{0,17,22\}$ & -- & -- & -- \\
$\mathbb{Z}_{9}\times\mathbb{Z}_{3}$ & C1, odd & $[[54,4,8]]$ & $\{(0,0),(0,1),(7,0)\}$ & $\{(0,0),(4,2),(7,0)\}$ & (7,0) & 9 & -- \\
$\mathbb{Z}_{9}\times\mathbb{Z}_{3}$ & D & $[[54,8,6]]$ & $\{(0,0),(1,0),(8,0)\}$ & $\{(0,0),(0,2),(3,1)\}$ & -- & -- & -- \\
$\mathbb{Z}_{9}\times\mathbb{Z}_{3}$ & D & $[[54,4,8]]$ & $\{(0,0),(3,1),(4,2)\}$ & $\{(0,0),(0,1),(2,2)\}$ & -- & -- & -- \\
$\mathbb{Z}_{27}$ & C1, odd & $[[54,4,8]]$ & $\{0,4,26\}$ & $\{0,16,26\}$ & 26 & 27 & -- \\
$\mathbb{Z}_{27}$ & D & $[[54,4,8]]$ & $\{0,19,23\}$ & $\{0,5,7\}$ & -- & -- & -- \\
\end{tabular}
\end{ruledtabular}
\end{table*}
\begin{table*}[!t]
\caption{\textbf{Pareto fronts, $N=36$ (exhaustive).}}
\label{tab:codes-b}
\footnotesize
\begin{ruledtabular}
\begin{tabular}{llllllcc}
group & class & $[[n,k,d]]$ & $A$ & $B$ & $\col$ & $\operatorname{ord}(\col)$ & $h(\col)$ \\
\colrule
$\mathbb{Z}_{2}\times\mathbb{Z}_{18}$ & C1, even & $[[72,8,6]]$ & $\{(0,0),(0,2),(0,7)\}$ & $\{(0,0),(0,1),(0,14)\}$ & (0,5) & 18 & 0 \\
$\mathbb{Z}_{2}\times\mathbb{Z}_{18}$ & C1, even & $[[72,4,8]]$ & $\{(0,0),(0,13),(1,5)\}$ & $\{(0,0),(0,13),(1,2)\}$ & (0,13) & 18 & 0 \\
$\mathbb{Z}_{2}\times\mathbb{Z}_{18}$ & C1, odd & $[[72,12,4]]$ & $\{(0,0),(1,6),(1,12)\}$ & $\{(0,0),(0,3),(0,6)\}$ & (0,6) & 3 & -- \\
$\mathbb{Z}_{2}\times\mathbb{Z}_{18}$ & C1, odd & $[[72,8,6]]$ & $\{(0,0),(0,5),(0,10)\}$ & $\{(0,0),(0,1),(0,8)\}$ & (0,10) & 9 & -- \\
$\mathbb{Z}_{2}\times\mathbb{Z}_{18}$ & C1, odd & $[[72,4,8]]$ & $\{(0,0),(0,10),(0,17)\}$ & $\{(0,0),(0,10),(1,5)\}$ & (0,10) & 9 & -- \\
$\mathbb{Z}_{2}\times\mathbb{Z}_{18}$ & D & $[[72,8,6]]$ & $\{(0,0),(0,11),(1,4)\}$ & $\{(0,0),(0,17),(1,16)\}$ & -- & -- & -- \\
$\mathbb{Z}_{2}\times\mathbb{Z}_{18}$ & D & $[[72,4,8]]$ & $\{(0,0),(1,14),(1,16)\}$ & $\{(0,0),(1,1),(1,11)\}$ & -- & -- & -- \\
$\mathbb{Z}_{3}\times\mathbb{Z}_{12}$ & C1, even & $[[72,12,4]]$ & $\{(0,0),(0,5),(0,7)\}$ & $\{(0,0),(0,10),(0,11)\}$ & (0,2) & 6 & 1 \\
$\mathbb{Z}_{3}\times\mathbb{Z}_{12}$ & C1, even & $[[72,8,6]]$ & $\{(0,0),(0,2),(2,6)\}$ & $\{(0,0),(2,0),(2,10)\}$ & (0,2) & 6 & 1 \\
$\mathbb{Z}_{3}\times\mathbb{Z}_{12}$ & C1, even & $[[72,4,8]]$ & $\{(0,0),(1,3),(2,0)\}$ & $\{(0,0),(1,10),(2,1)\}$ & (1,3) & 12 & 0 \\
$\mathbb{Z}_{3}\times\mathbb{Z}_{12}$ & C1, odd & $[[72,12,4]]$ & $\{(0,0),(1,6),(2,6)\}$ & $\{(0,0),(1,9),(2,9)\}$ & (1,0) & 3 & -- \\
$\mathbb{Z}_{3}\times\mathbb{Z}_{12}$ & C1, odd & $[[72,4,6]]$ & $\{(0,0),(0,8),(2,1)\}$ & $\{(0,0),(2,7),(2,11)\}$ & (0,8) & 3 & -- \\
$\mathbb{Z}_{3}\times\mathbb{Z}_{12}$ & D & $[[72,16,4]]$ & $\{(0,0),(1,10),(2,2)\}$ & $\{(0,0),(0,8),(0,10)\}$ & -- & -- & -- \\
$\mathbb{Z}_{3}\times\mathbb{Z}_{12}$ & D & $[[72,8,8]]$ & $\{(0,0),(0,1),(0,8)\}$ & $\{(0,0),(1,8),(2,7)\}$ & -- & -- & -- \\
$\mathbb{Z}_{6}\times\mathbb{Z}_{6}$ & C1, even & $[[72,8,6]]$ & $\{(0,0),(4,3),(5,3)\}$ & $\{(0,0),(1,5),(4,3)\}$ & (4,3) & 6 & 0 \\
$\mathbb{Z}_{6}\times\mathbb{Z}_{6}$ & C1, even & $[[72,4,8]]$ & $\{(0,0),(0,5),(4,4)\}$ & $\{(0,0),(1,1),(1,2)\}$ & (0,5) & 6 & 0 \\
$\mathbb{Z}_{6}\times\mathbb{Z}_{6}$ & C1, odd & $[[72,12,4]]$ & $\{(0,0),(3,2),(3,4)\}$ & $\{(0,0),(3,1),(3,5)\}$ & (0,2) & 3 & -- \\
$\mathbb{Z}_{6}\times\mathbb{Z}_{6}$ & C1, odd & $[[72,4,8]]$ & $\{(0,0),(4,3),(4,4)\}$ & $\{(0,0),(1,3),(3,5)\}$ & (4,4) & 3 & -- \\
$\mathbb{Z}_{6}\times\mathbb{Z}_{6}$ & D & $[[72,16,4]]$ & $\{(0,0),(0,2),(3,4)\}$ & $\{(0,0),(1,0),(5,0)\}$ & -- & -- & -- \\
$\mathbb{Z}_{6}\times\mathbb{Z}_{6}$ & D & $[[72,12,6]]$ & $\{(0,0),(0,1),(3,5)\}$ & $\{(0,0),(1,0),(2,3)\}$ & -- & -- & -- \\
$\mathbb{Z}_{6}\times\mathbb{Z}_{6}$ & D & $[[72,4,8]]$ & $\{(0,0),(2,0),(4,5)\}$ & $\{(0,0),(1,3),(2,4)\}$ & -- & -- & -- \\
$\mathbb{Z}_{36}$ & C1, even & $[[72,12,4]]$ & $\{0,15,21\}$ & $\{0,3,33\}$ & 6 & 6 & 1 \\
$\mathbb{Z}_{36}$ & C1, even & $[[72,8,6]]$ & $\{0,2,28\}$ & $\{0,10,32\}$ & 26 & 18 & 1 \\
$\mathbb{Z}_{36}$ & C1, even & $[[72,4,8]]$ & $\{0,22,23\}$ & $\{0,13,29\}$ & 23 & 36 & 0 \\
$\mathbb{Z}_{36}$ & C1, odd & $[[72,12,4]]$ & $\{0,12,15\}$ & $\{0,6,12\}$ & 12 & 3 & -- \\
$\mathbb{Z}_{36}$ & C1, odd & $[[72,8,6]]$ & $\{0,8,22\}$ & $\{0,8,10\}$ & 8 & 9 & -- \\
$\mathbb{Z}_{36}$ & C1, odd & $[[72,4,8]]$ & $\{0,23,31\}$ & $\{0,7,8\}$ & 8 & 9 & -- \\
$\mathbb{Z}_{36}$ & D & $[[72,8,6]]$ & $\{0,16,20\}$ & $\{0,14,28\}$ & -- & -- & -- \\
$\mathbb{Z}_{36}$ & D & $[[72,4,\le10]]$ & $\{0,20,25\}$ & $\{0,1,29\}$ & -- & -- & -- \\
\end{tabular}
\end{ruledtabular}
\end{table*}
\begin{table*}[!t]
\caption{\textbf{Pareto fronts, $N\in\{45,54,72\}$ (sampled).}}
\label{tab:codes-c}
\footnotesize
\begin{ruledtabular}
\begin{tabular}{llllllcc}
group & class & $[[n,k,d]]$ & $A$ & $B$ & $\col$ & $\operatorname{ord}(\col)$ & $h(\col)$ \\
\colrule
$\mathbb{Z}_{15}\times\mathbb{Z}_{3}$ & C1, odd & $[[90,12,6]]$ & $\{(0,0),(8,2),(13,1)\}$ & $\{(0,0),(1,1),(14,2)\}$ & (13,1) & 15 & -- \\
$\mathbb{Z}_{15}\times\mathbb{Z}_{3}$ & C1, odd & $[[90,4,\le10]]$ & $\{(0,0),(1,1),(11,2)\}$ & $\{(0,0),(6,1),(10,2)\}$ & (11,2) & 15 & -- \\
$\mathbb{Z}_{15}\times\mathbb{Z}_{3}$ & D & $[[90,24,4]]$ & $\{(0,0),(7,0),(9,0)\}$ & $\{(0,0),(11,0),(12,0)\}$ & -- & -- & -- \\
$\mathbb{Z}_{15}\times\mathbb{Z}_{3}$ & D & $[[90,12,6]]$ & $\{(0,0),(2,0),(13,1)\}$ & $\{(0,0),(7,1),(8,0)\}$ & -- & -- & -- \\
$\mathbb{Z}_{15}\times\mathbb{Z}_{3}$ & D & $[[90,8,\le10]]$ & $\{(0,0),(10,0),(14,0)\}$ & $\{(0,0),(8,1),(13,2)\}$ & -- & -- & -- \\
$\mathbb{Z}_{45}$ & C1, odd & $[[90,12,6]]$ & $\{0,12,33\}$ & $\{0,24,39\}$ & 21 & 15 & -- \\
$\mathbb{Z}_{45}$ & C1, odd & $[[90,4,\le10]]$ & $\{0,25,41\}$ & $\{0,25,38\}$ & 25 & 9 & -- \\
$\mathbb{Z}_{45}$ & D & $[[90,24,4]]$ & $\{0,6,27\}$ & $\{0,9,12\}$ & -- & -- & -- \\
$\mathbb{Z}_{45}$ & D & $[[90,8,\le8]]$ & $\{0,11,29\}$ & $\{0,2,39\}$ & -- & -- & -- \\
$\mathbb{Z}_{45}$ & D & $[[90,4,\le12]]$ & $\{0,14,19\}$ & $\{0,1,17\}$ & -- & -- & -- \\
$\mathbb{Z}_{9}\times\mathbb{Z}_{6}$ & C1, even & $[[108,12,6]]$ & $\{(0,0),(4,4),(5,5)\}$ & $\{(0,0),(7,1),(8,2)\}$ & (1,1) & 18 & 0 \\
$\mathbb{Z}_{9}\times\mathbb{Z}_{6}$ & C1, even & $[[108,4,\le12]]$ & $\{(0,0),(5,4),(7,1)\}$ & $\{(0,0),(2,0),(7,3)\}$ & (2,3) & 18 & 0 \\
$\mathbb{Z}_{9}\times\mathbb{Z}_{6}$ & C1, odd & $[[108,12,6]]$ & $\{(0,0),(1,1),(2,2)\}$ & $\{(0,0),(2,5),(7,4)\}$ & (2,2) & 9 & -- \\
$\mathbb{Z}_{9}\times\mathbb{Z}_{6}$ & C1, odd & $[[108,8,\le8]]$ & $\{(0,0),(1,2),(4,0)\}$ & $\{(0,0),(0,2),(4,2)\}$ & (4,0) & 9 & -- \\
$\mathbb{Z}_{9}\times\mathbb{Z}_{6}$ & C1, odd & $[[108,4,\le12]]$ & $\{(0,0),(4,5),(5,4)\}$ & $\{(0,0),(1,0),(5,4)\}$ & (5,4) & 9 & -- \\
$\mathbb{Z}_{9}\times\mathbb{Z}_{6}$ & D & $[[108,24,4]]$ & $\{(0,0),(0,1),(0,2)\}$ & $\{(0,0),(3,5),(6,4)\}$ & -- & -- & -- \\
$\mathbb{Z}_{9}\times\mathbb{Z}_{6}$ & D & $[[108,16,6]]$ & $\{(0,0),(6,2),(6,4)\}$ & $\{(0,0),(2,2),(4,4)\}$ & -- & -- & -- \\
$\mathbb{Z}_{9}\times\mathbb{Z}_{6}$ & D & $[[108,8,\le10]]$ & $\{(0,0),(3,1),(3,5)\}$ & $\{(0,0),(2,3),(7,0)\}$ & -- & -- & -- \\
$\mathbb{Z}_{9}\times\mathbb{Z}_{6}$ & D & $[[108,4,\le12]]$ & $\{(0,0),(2,2),(4,3)\}$ & $\{(0,0),(1,4),(5,4)\}$ & -- & -- & -- \\
$\mathbb{Z}_{18}\times\mathbb{Z}_{3}$ & C1, even & $[[108,12,6]]$ & $\{(0,0),(10,1),(17,2)\}$ & $\{(0,0),(11,2),(16,1)\}$ & (7,1) & 18 & 0 \\
$\mathbb{Z}_{18}\times\mathbb{Z}_{3}$ & C1, even & $[[108,4,\le12]]$ & $\{(0,0),(10,2),(11,1)\}$ & $\{(0,0),(4,1),(17,1)\}$ & (1,2) & 18 & 0 \\
$\mathbb{Z}_{18}\times\mathbb{Z}_{3}$ & C1, odd & $[[108,12,6]]$ & $\{(0,0),(8,0),(13,0)\}$ & $\{(0,0),(7,0),(17,0)\}$ & (8,0) & 9 & -- \\
$\mathbb{Z}_{18}\times\mathbb{Z}_{3}$ & C1, odd & $[[108,8,\le8]]$ & $\{(0,0),(12,2),(16,1)\}$ & $\{(0,0),(0,1),(2,2)\}$ & (16,1) & 9 & -- \\
$\mathbb{Z}_{18}\times\mathbb{Z}_{3}$ & C1, odd & $[[108,4,\le12]]$ & $\{(0,0),(0,2),(16,2)\}$ & $\{(0,0),(13,2),(16,0)\}$ & (16,0) & 9 & -- \\
$\mathbb{Z}_{18}\times\mathbb{Z}_{3}$ & D & $[[108,24,4]]$ & $\{(0,0),(3,2),(6,1)\}$ & $\{(0,0),(3,0),(15,0)\}$ & -- & -- & -- \\
$\mathbb{Z}_{18}\times\mathbb{Z}_{3}$ & D & $[[108,16,6]]$ & $\{(0,0),(6,1),(6,2)\}$ & $\{(0,0),(8,1),(16,2)\}$ & -- & -- & -- \\
$\mathbb{Z}_{18}\times\mathbb{Z}_{3}$ & D & $[[108,8,\le10]]$ & $\{(0,0),(7,2),(17,1)\}$ & $\{(0,0),(0,1),(15,2)\}$ & -- & -- & -- \\
$\mathbb{Z}_{18}\times\mathbb{Z}_{3}$ & D & $[[108,4,\le12]]$ & $\{(0,0),(1,1),(17,1)\}$ & $\{(0,0),(4,1),(11,2)\}$ & -- & -- & -- \\
$\mathbb{Z}_{54}$ & C1, even & $[[108,12,6]]$ & $\{0,6,39\}$ & $\{0,39,51\}$ & 39 & 18 & 0 \\
$\mathbb{Z}_{54}$ & C1, even & $[[108,4,\le12]]$ & $\{0,1,29\}$ & $\{0,37,38\}$ & 1 & 54 & 0 \\
$\mathbb{Z}_{54}$ & C1, odd & $[[108,12,6]]$ & $\{0,3,33\}$ & $\{0,30,42\}$ & 30 & 9 & -- \\
$\mathbb{Z}_{54}$ & C1, odd & $[[108,8,\le8]]$ & $\{0,2,22\}$ & $\{0,22,50\}$ & 22 & 27 & -- \\
$\mathbb{Z}_{54}$ & C1, odd & $[[108,4,\le12]]$ & $\{0,19,32\}$ & $\{0,32,40\}$ & 32 & 27 & -- \\
$\mathbb{Z}_{54}$ & D & $[[108,12,6]]$ & $\{0,6,30\}$ & $\{0,21,42\}$ & -- & -- & -- \\
$\mathbb{Z}_{54}$ & D & $[[108,8,\le8]]$ & $\{0,2,4\}$ & $\{0,14,22\}$ & -- & -- & -- \\
$\mathbb{Z}_{54}$ & D & $[[108,4,\le12]]$ & $\{0,34,53\}$ & $\{0,10,38\}$ & -- & -- & -- \\
$\mathbb{Z}_{12}\times\mathbb{Z}_{6}$ & C1, even & $[[144,24,4]]$ & $\{(0,0),(7,0),(11,3)\}$ & $\{(0,0),(7,3),(8,3)\}$ & (4,3) & 6 & 0 \\
$\mathbb{Z}_{12}\times\mathbb{Z}_{6}$ & C1, even & $[[144,16,6]]$ & $\{(0,0),(10,0),(10,4)\}$ & $\{(0,0),(2,0),(8,2)\}$ & (10,0) & 6 & 1 \\
$\mathbb{Z}_{12}\times\mathbb{Z}_{6}$ & C1, even & $[[144,8,\le10]]$ & $\{(0,0),(1,0),(7,5)\}$ & $\{(0,0),(4,3),(9,4)\}$ & (7,5) & 12 & 0 \\
$\mathbb{Z}_{12}\times\mathbb{Z}_{6}$ & C1, even & $[[144,4,\le12]]$ & $\{(0,0),(7,1),(11,2)\}$ & $\{(0,0),(5,3),(7,1)\}$ & (7,1) & 12 & 0 \\
$\mathbb{Z}_{12}\times\mathbb{Z}_{6}$ & C1, odd & $[[144,24,4]]$ & $\{(0,0),(0,2),(6,1)\}$ & $\{(0,0),(0,1),(0,2)\}$ & (0,2) & 3 & -- \\
$\mathbb{Z}_{12}\times\mathbb{Z}_{6}$ & C1, odd & $[[144,8,\le8]]$ & $\{(0,0),(6,2),(10,4)\}$ & $\{(0,0),(0,1),(8,5)\}$ & (4,2) & 3 & -- \\
$\mathbb{Z}_{12}\times\mathbb{Z}_{6}$ & D & $[[144,32,4]]$ & $\{(0,0),(8,2),(10,4)\}$ & $\{(0,0),(2,0),(10,0)\}$ & -- & -- & -- \\
$\mathbb{Z}_{12}\times\mathbb{Z}_{6}$ & D & $[[144,24,6]]$ & $\{(0,0),(2,4),(10,5)\}$ & $\{(0,0),(0,5),(6,1)\}$ & -- & -- & -- \\
$\mathbb{Z}_{12}\times\mathbb{Z}_{6}$ & D & $[[144,16,\le8]]$ & $\{(0,0),(3,2),(3,4)\}$ & $\{(0,0),(4,2),(11,4)\}$ & -- & -- & -- \\
$\mathbb{Z}_{12}\times\mathbb{Z}_{6}$ & D & $[[144,12,\le12]]$ & $\{(0,0),(7,4),(11,5)\}$ & $\{(0,0),(2,1),(7,2)\}$ & -- & -- & -- \\
$\mathbb{Z}_{12}\times\mathbb{Z}_{6}$ & D & $[[144,4,\le14]]$ & $\{(0,0),(5,5),(8,1)\}$ & $\{(0,0),(1,4),(7,5)\}$ & -- & -- & -- \\
$\mathbb{Z}_{72}$ & C1, even & $[[144,24,4]]$ & $\{0,60,66\}$ & $\{0,30,42\}$ & 60 & 6 & 2 \\
$\mathbb{Z}_{72}$ & C1, even & $[[144,16,6]]$ & $\{0,40,68\}$ & $\{0,8,52\}$ & 28 & 18 & 2 \\
$\mathbb{Z}_{72}$ & C1, even & $[[144,12,\le8]]$ & $\{0,3,51\}$ & $\{0,51,66\}$ & 51 & 24 & 0 \\
$\mathbb{Z}_{72}$ & C1, even & $[[144,4,\le14]]$ & $\{0,29,64\}$ & $\{0,17,37\}$ & 35 & 72 & 0 \\
$\mathbb{Z}_{72}$ & C1, odd & $[[144,24,4]]$ & $\{0,42,66\}$ & $\{0,48,60\}$ & 24 & 3 & -- \\
$\mathbb{Z}_{72}$ & C1, odd & $[[144,16,6]]$ & $\{0,4,8\}$ & $\{0,44,64\}$ & 8 & 9 & -- \\
$\mathbb{Z}_{72}$ & C1, odd & $[[144,8,\le8]]$ & $\{0,28,68\}$ & $\{0,14,46\}$ & 40 & 9 & -- \\
$\mathbb{Z}_{72}$ & C1, odd & $[[144,4,\le14]]$ & $\{0,10,50\}$ & $\{0,40,65\}$ & 40 & 9 & -- \\
$\mathbb{Z}_{72}$ & D & $[[144,16,6]]$ & $\{0,32,52\}$ & $\{0,56,64\}$ & -- & -- & -- \\
$\mathbb{Z}_{72}$ & D & $[[144,12,\le8]]$ & $\{0,42,48\}$ & $\{0,39,60\}$ & -- & -- & -- \\
$\mathbb{Z}_{72}$ & D & $[[144,8,\le10]]$ & $\{0,38,58\}$ & $\{0,44,70\}$ & -- & -- & -- \\
$\mathbb{Z}_{72}$ & D & $[[144,4,\le16]]$ & $\{0,25,62\}$ & $\{0,8,34\}$ & -- & -- & -- \\
\end{tabular}
\end{ruledtabular}
\end{table*}

\paragraph*{Explicit depth-six schedules.}
Table~\ref{tab:template72} of the main text gives the class schedule of the
$[[72,8,6]]$ code. Table~\ref{tab:templates-more} gives class schedules,
found by SAT on the class system and verified by $\F$ simulation of the
pullback, for the $[[108,12,6]]$ code over $\mathbb{Z}_{54}$, the
$[[108,4,\le12]]$ code over $\mathbb{Z}_9\times\mathbb{Z}_6$, the
$[[144,12,\le8]]$ code over $\mathbb{Z}_{72}$ and the $[[144,8,\le10]]$
code over $\mathbb{Z}_{12}\times\mathbb{Z}_6$, all with $h(\col)=0$ and
$q=2$, and Table~\ref{tab:template144-16} the $q=4$ class schedule of the
$[[144,16,6]]$ code over $\mathbb{Z}_{12}\times\mathbb{Z}_6$, whose
collision element $(10,0)$ has $2$-height $1$. The $[[144,16,6]]$ code
over $\mathbb{Z}_{72}$ ($A=\{0,40,68\}$, $B=\{0,8,52\}$, $\col=28$ of
$2$-height $2$) was scheduled at $q=8$ in the same way; its key
$(0,0,4;0,0,4)$ generates the subgroup $\{0,4\}$ of $\mathbb{Z}_8$, so its
class system is four copies of a $q=2$ system, as described in Methods.
Information-set decoding with $10^4$ iterations did not lower any of the
distance bounds of these codes.

\begin{table*}[!t]
\caption{\textbf{Class schedules ($q=2$) of four depth-six codes.} Layers
by class $c$ and edge type ($XL$/$XR$: $X$-check $L$-/$R$-edges, $ZL$/$ZR$:
$Z$-check $L$-/$R$-edges), labels $1,2,3$ in order. Collisions are given
as label pairs $(s,t)$; $\varepsilon$ reduces the last coordinate modulo
$2$ ($g\bmod2$ on cyclic groups).}
\label{tab:templates-more}
\begin{ruledtabular}
\begin{tabular}{lcccc}
code, $c$ & $XL$ & $XR$ & $ZL$ & $ZR$ \\
\colrule
\multicolumn{5}{l}{$[[108,12,6]]$: $\mathbb{Z}_{54}$, $A=\{0,6,39\}$, $B=\{0,39,51\}$,}\\
\multicolumn{5}{l}{\quad collision $(1,2),(3,1)$, key $(0,0,1;0,1,1)$}\\
$c=0$ & $(3,4,1)$ & $(6,2,5)$ & $(1,4,3)$ & $(2,5,6)$ \\
$c=1$ & $(6,2,5)$ & $(3,1,4)$ & $(5,6,2)$ & $(1,4,3)$ \\
\colrule
\multicolumn{5}{l}{$[[108,4,\le12]]$: $\mathbb{Z}_9\times\mathbb{Z}_6$, $A=\{(0,0),(5,4),(7,1)\}$,}\\
\multicolumn{5}{l}{\quad $B=\{(0,0),(2,0),(7,3)\}$, collision $(2,1),(3,3)$, key $(0,0,1;0,0,1)$}\\
$c=0$ & $(4,2,5)$ & $(3,6,1)$ & $(6,1,4)$ & $(5,2,3)$ \\
$c=1$ & $(1,6,3)$ & $(2,5,4)$ & $(3,2,5)$ & $(6,4,1)$ \\
\colrule
\multicolumn{5}{l}{$[[144,12,\le8]]$: $\mathbb{Z}_{72}$, $A=\{0,3,51\}$, $B=\{0,51,66\}$,}\\
\multicolumn{5}{l}{\quad collision $(1,2),(3,1)$, key $(0,1,1;0,1,0)$}\\
$c=0$ & $(4,2,5)$ & $(1,3,6)$ & $(5,3,2)$ & $(4,6,1)$ \\
$c=1$ & $(6,1,3)$ & $(5,2,4)$ & $(4,6,1)$ & $(2,3,5)$ \\
\colrule
\multicolumn{5}{l}{$[[144,8,\le10]]$: $\mathbb{Z}_{12}\times\mathbb{Z}_6$, $A=\{(0,0),(1,0),(7,5)\}$,}\\
\multicolumn{5}{l}{\quad $B=\{(0,0),(4,3),(9,4)\}$, collision $(1,2),(3,3)$, key $(0,0,1;0,1,0)$}\\
$c=0$ & $(4,2,5)$ & $(6,3,1)$ & $(1,4,6)$ & $(3,5,2)$ \\
$c=1$ & $(6,1,3)$ & $(5,2,4)$ & $(2,5,3)$ & $(1,6,4)$ \\
\end{tabular}
\end{ruledtabular}
\end{table*}

\begin{table}[!htbp]
\caption{\textbf{Class schedule ($q=4$) of the $[[144,16,6]]$ code.}
$G=\mathbb{Z}_{12}\times\mathbb{Z}_6$, $A=\{(0,0),(10,0),(10,4)\}$,
$B=\{(0,0),(2,0),(8,2)\}$, collision $(1,1),(2,2)$, $\col=(10,0)$,
$\varepsilon(i,j)=i\bmod4$, key $(0,2,2;0,2,0)$, $\colbar=2$.}
\label{tab:template144-16}
\begin{ruledtabular}
\begin{tabular}{lcccc}
$c$ & $XL$ & $XR$ & $ZL$ & $ZR$ \\
\colrule
0 & $(4,1,6)$ & $(2,5,3)$ & $(1,4,6)$ & $(5,2,3)$ \\
1 & $(4,5,2)$ & $(6,1,3)$ & $(5,3,2)$ & $(1,4,6)$ \\
2 & $(2,5,3)$ & $(6,4,1)$ & $(5,2,3)$ & $(4,6,1)$ \\
3 & $(6,3,1)$ & $(2,4,5)$ & $(1,6,4)$ & $(3,2,5)$ \\
\end{tabular}
\end{ruledtabular}
\end{table}

\paragraph*{Odd-order instances decided directly.}
Depth $6$ was decided by SAT for the best odd-order class-C1 codes of
Tables~\ref{tab:codes-b} and~\ref{tab:codes-c}: the $[[72,8,6]]$ code over
$\mathbb{Z}_{36}$ ($A=\{0,8,22\}$, $B=\{0,8,10\}$, $\col=8$ of order $9$)
is satisfiable ($60$~s), as is the $[[72,12,4]]$ code over
$\mathbb{Z}_6\times\mathbb{Z}_6$ ($A=\{(0,0),(3,2),(3,4)\}$,
$B=\{(0,0),(3,1),(3,5)\}$, $\col=(0,2)$ of order $3$; $7$~s), both
schedules verified by $\F$ simulation; the $[[90,12,6]]$ code over
$\mathbb{Z}_{15}\times\mathbb{Z}_3$ ($A=\{(0,0),(8,2),(13,1)\}$,
$B=\{(0,0),(1,1),(14,2)\}$, $\col=(13,1)$ of order $15$) is unsatisfiable:
first by \textsc{CaDiCaL} in $523$~s without proof logging, then re-decided
by \textsc{kissat} with DRAT logging in $600$~s, giving a $2.8$~GB proof
verified by \textsc{drat-trim}. The $[[90,12,6]]$ code over
$\mathbb{Z}_{45}$ ($A=\{0,12,33\}$, $B=\{0,24,39\}$, $\col=21$ of order
$15$), whose earlier PySAT run crashed and which was therefore left
undecided, is also unsatisfiable: \textsc{kissat} refutes it in $967$~s and
\textsc{drat-trim} verifies the $2.4$~GB proof. Both certificates are in
Table~\ref{tab:certs-full} and Fig.~\ref{fig:certs}.

\FloatBarrier

\end{document}